% Standalone build:
%   latexmk -pdf -interaction=nonstopmode -halt-on-error
%     crouzeix_semester_workshop.tex
% latexmk performs the multiple passes required by cross-references and
% the inline bibliography; the completed build has a warning-free log.
\documentclass[11pt,oneside,openany]{book}

\usepackage[T1]{fontenc}
\usepackage[utf8]{inputenc}
\usepackage[letterpaper,margin=1in,headheight=15pt]{geometry}
\usepackage{lmodern}
\usepackage{amsmath,amssymb,amsthm,mathtools}
\usepackage{bm}
\usepackage{booktabs,longtable,tabularx,array}
\usepackage{enumitem}
\usepackage[dvipsnames]{xcolor}
\usepackage{microtype}
\usepackage{fancyhdr}
\usepackage{tikz}
\usetikzlibrary{arrows.meta,positioning,calc,shapes.geometric}
\usepackage[most]{tcolorbox}
\usepackage{hyperref}
\usepackage[nameinlink,capitalise,noabbrev]{cleveref}

\definecolor{WorkshopBlue}{HTML}{214E73}
\definecolor{WorkshopTeal}{HTML}{197278}
\definecolor{WorkshopGold}{HTML}{C98B2E}
\definecolor{WorkshopRed}{HTML}{A23E48}
\definecolor{WorkshopInk}{HTML}{23313F}
\definecolor{WorkshopPale}{HTML}{F4F7FA}

\hypersetup{
  colorlinks=true,
  linkcolor=WorkshopBlue,
  citecolor=WorkshopTeal,
  urlcolor=WorkshopRed,
  pdftitle={The Numerical Range Is a 2-Spectral Set: A Semester Workshop},
  pdfauthor={Workshop notes based on the manuscript by Shanmu Jin}
}

\pagestyle{fancy}
\fancyhf{}
\fancyhead[L]{\small\itshape The numerical range is a \(2\)-spectral set}
\fancyhead[R]{\small\itshape Semester workshop}
\fancyfoot[C]{\thepage}
\renewcommand{\headrulewidth}{0.4pt}

\setcounter{secnumdepth}{3}
\setcounter{tocdepth}{2}
\setlist[itemize]{topsep=4pt,itemsep=2pt}
\setlist[enumerate]{topsep=4pt,itemsep=3pt}
\allowdisplaybreaks
\newcolumntype{L}[1]{>{\raggedright\arraybackslash}p{#1}}

\newtheorem{theorem}{Theorem}[chapter]
\newtheorem{lemma}[theorem]{Lemma}
\newtheorem{proposition}[theorem]{Proposition}
\newtheorem{corollary}[theorem]{Corollary}

\theoremstyle{definition}
\newtheorem{definition}[theorem]{Definition}
\newtheorem{example}[theorem]{Example}
\newtheorem{exercise}{Exercise}[chapter]

\theoremstyle{remark}
\newtheorem{remark}[theorem]{Remark}

\crefname{theorem}{theorem}{theorems}
\crefname{lemma}{lemma}{lemmas}
\crefname{proposition}{proposition}{propositions}
\crefname{corollary}{corollary}{corollaries}
\crefname{definition}{definition}{definitions}
\crefname{example}{example}{examples}
\crefname{exercise}{exercise}{exercises}

\newtcolorbox{learninggoals}{
  enhanced,breakable,colback=WorkshopPale,colframe=WorkshopBlue,
  title={Learning goals},fonttitle=\bfseries,boxrule=0.7pt,arc=2pt}
\newtcolorbox{workshopplan}{
  enhanced,breakable,colback=white,colframe=WorkshopTeal,
  title={Workshop plan},fonttitle=\bfseries,boxrule=0.7pt,arc=2pt}
\newtcolorbox{intuitionbox}[1][Intuition]{
  enhanced,breakable,colback=WorkshopTeal!5,colframe=WorkshopTeal,
  title={#1},fonttitle=\bfseries,boxrule=0.7pt,arc=2pt}
\newtcolorbox{checkpoint}[1][Checkpoint]{
  enhanced,breakable,colback=WorkshopGold!8,colframe=WorkshopGold!90!black,
  title={#1},fonttitle=\bfseries,boxrule=0.7pt,arc=2pt}
\newtcolorbox{pitfall}[1][Common pitfall]{
  enhanced,breakable,colback=WorkshopRed!5,colframe=WorkshopRed,
  title={#1},fonttitle=\bfseries,boxrule=0.7pt,arc=2pt}
\newtcolorbox{proofmap}[1][Proof map]{
  enhanced,breakable,colback=WorkshopBlue!4,colframe=WorkshopBlue,
  title={#1},fonttitle=\bfseries,boxrule=0.9pt,arc=2pt}
\newtcolorbox{blackbox}[1][Classical input]{
  enhanced,breakable,colback=black!2,colframe=black!55,
  title={#1},fonttitle=\bfseries,boxrule=0.6pt,arc=2pt}

\newcommand{\C}{\mathbb C}
\newcommand{\R}{\mathbb R}
\newcommand{\D}{\mathbb D}
\newcommand{\Torus}{\mathbb T}
\newcommand{\Mn}{\mathbb M_n(\C)}
\newcommand{\Mm}{\mathbb M_m(\C)}
\newcommand{\alg}{\operatorname{alg}}
\newcommand{\diag}{\operatorname{diag}}
\newcommand{\dist}{\operatorname{dist}}
\newcommand{\conv}{\operatorname{conv}}
\newcommand{\RePart}{\operatorname{Re}}
\newcommand{\ImPart}{\operatorname{Im}}
\newcommand{\spec}{\sigma}
\newcommand{\rank}{\operatorname{rank}}
\newcommand{\Span}{\operatorname{span}}
\newcommand{\ip}[2]{\left\langle #1,#2\right\rangle}
\newcommand{\norm}[1]{\left\lVert #1\right\rVert}
\newcommand{\abs}[1]{\left\lvert #1\right\rvert}
\newcommand{\set}[1]{\left\{#1\right\}}
\newcommand{\eps}{\varepsilon}
\newcommand{\olD}{\overline{\D}}
\newcommand{\WA}{W(A)}

\DeclareMathOperator{\supp}{supp}

\title{\bfseries The Numerical Range Is a \(2\)-Spectral Set\\[0.4em]
\Large A Semester Workshop from First Principles}
\author{Pedagogical notes based on the manuscript\\
\emph{The numerical range is a \(2\)-spectral set} by Shanmu Jin}
\date{Workshop edition, July 2026}

\begin{document}

\frontmatter
\hypersetup{pageanchor=false}
\maketitle
\hypersetup{pageanchor=true}

\chapter*{Preface}
\addcontentsline{toc}{chapter}{Preface}

These notes turn the proof in the source manuscript
\cite{JinManuscript} into a semester-long workshop.  The intended reader
does \emph{not} need prior exposure to Crouzeix's conjecture, numerical
ranges, spectral sets, positive kernels, double-layer potentials, or
matrix-valued Herglotz theory.  The phrase ``from ABC/101'' here means
that no problem-specific vocabulary or folklore is assumed.  A first
course in finite-dimensional linear algebra and the basic definitions
of complex analysis are useful.  The notes review the needed objects
from the ground up and isolate deeper classical results---such as
Schur triangularization, Cauchy's formula, and Runge approximation---in
clearly marked boxes with proof blueprints.

The central theorem is easy to state.  For a complex square matrix
\(A\), define its numerical range by
\[
  W(A)=\set{x^*Ax:\ x\in\C^n,\ \norm{x}_2=1}.
\]
For every complex polynomial \(p\), the source manuscript presents a
proof of
\[
  \boxed{\quad
  \norm{p(A)}
  \leq 2\max_{z\in W(A)}\abs{p(z)}.
  \quad}
\]
The Week 3 nilpotent example shows that no valid universal constant
can be smaller than \(2\); thus the manuscript's claimed bound would
be sharp.

What makes the proof worth a workshop is not only the conclusion.
The new idea is a reusable two-stage mechanism:
\[
  \text{structured cancellation in a positive kernel}
  \quad\Longrightarrow\quad
  \text{rigidity of two weighted Gramians}.
\]
An unknown analytic correction is not bounded or discarded; a carefully
chosen sample at the origin cancels it \emph{exactly}.  The surviving
inequality is converted into the sharp constant by a Stein identity.
The middle third of the course studies this mechanism line by line.

\begin{pitfall}[Status and purpose of these notes]
These notes are a pedagogical expansion of the stated source
manuscript, which the repository describes as a candidate proof with
formal peer review pending.  They are not an independent refereeing
report.  They expose every logical dependency so that a workshop can
check the argument rather than merely quote it.  When a classical
theorem such as Runge's theorem or the scalar Herglotz representation
is used, its precise statement, role, and proof blueprint are supplied.
\end{pitfall}

\section*{How to use the material}

Each numbered chapter is one workshop week.  A typical meeting can use
two hours for guided theory and one hour for problems.  The boxes have
consistent roles:

\begin{itemize}
  \item \textbf{Learning goals} say what a participant should be able
  to do by the end of the week.
  \item \textbf{Intuition} translates a formal step into a mental model.
  \item \textbf{Checkpoint} pauses the proof and records exactly what
  has been earned.
  \item \textbf{Common pitfall} marks a tempting but invalid shortcut.
  \item \textbf{Classical input} isolates a theorem whose full proof
  belongs to a standard prerequisite course.
\end{itemize}

Exercises are marked \emph{warm-up}, \emph{core}, or \emph{stretch}.
Hints and selected solutions appear in \cref{app:solutions}.  A reader
following only the core route should complete all core exercises and
the proof-reconstruction checklists in Weeks 12 and 14.

\begin{intuitionbox}[A short Week 0 diagnostic]
Before Week 1, check that you can multiply \(2\times2\) complex
matrices, solve a two-equation linear system, and differentiate
\(z^k\).  Everything else needed for the argument is either developed
in the notes or labeled as a classical input.  If the words
``spectrum,'' ``holomorphic,'' or ``positive semidefinite'' are new,
that is expected: Weeks 1, 2, and 5 introduce them.
\end{intuitionbox}

\tableofcontents

\chapter*{Semester road map}
\addcontentsline{toc}{chapter}{Semester road map}

\begin{longtable}{@{}L{0.07\textwidth}L{0.25\textwidth}
                         L{0.42\textwidth}L{0.17\textwidth}@{}}
\toprule
\textbf{Week}&\textbf{Theme}&\textbf{Proof milestone}&\textbf{Deliverable}\\
\midrule
\endhead
1&Matrix language&Norms, adjoints, positivity, and congruence&
Fluency sheet\\
2&Spectra and matrix functions&Generated algebras and interpolation&
Functional-calculus map\\
3&Numerical ranges&Geometry, convexity, and sharp \(2\times2\) example&
Geometry problem set\\
4&Spectral sets&The conjecture and normalization&
One-page problem brief\\
5&Complex and convex analysis&Resolvents, Cauchy calculus, smooth boundaries&
Contour-calculus sheet\\
6&Positive maps and measures&The language behind the boundary calculus&
Positivity audit\\
7&Herglotz kernels&Positive-real functions become matrix inequalities&
Kernel sampling lab\\
8&Double-layer calculus&Construct the unital positive map \(\Phi\)&
Derivation portfolio\\
9&Cayley completion&Produce \(H\) and locate its resolvent defect&
Completion proof\\
10&Diagonal algebra and cancellation&Diagonalize the defect and cancel it
with half-eigenvalue and origin samples&Change-of-basis audit\\
11&Surviving inequality&Compute the known blocks and identify two Gramians&
Cancellation calculation\\
12&Gramian rigidity&Two scales, an eigenvector test, and a Stein identity&
Core theorem oral proof\\
13&Removing hypotheses&Three approximation layers&
Limit ledger\\
14&Main theorem&Assemble the polynomial and rational conclusions&
Full proof write-up\\
15&Consequences and computation&Holomorphic calculus, compression, algorithms&
Numerical experiment\\
16&New-tool clinic&General scale, adjoint alignment, and open directions&
Capstone presentation\\
\bottomrule
\end{longtable}

\begin{proofmap}[The whole proof in one page]
\begin{center}
\begin{tikzpicture}[
  node distance=8mm and 10mm,
  box/.style={draw=WorkshopBlue,rounded corners,align=center,
              fill=WorkshopBlue!4,text width=3.3cm,minimum height=1.05cm},
  arr/.style={-{Stealth[length=2.2mm]},thick,WorkshopTeal}
]
\node[box] (geom) {Numerical range inside a smooth convex domain};
\node[box,right=of geom] (phi) {Positive double-layer map \(\Phi\)};
\node[box,right=of phi] (cayley) {Cayley family \(H(w)=\Phi((1+wf)/(1-wf))\)};
\node[box,below=of cayley] (alg) {Resolvent defect in \(\alg(T^*)\)};
\node[box,left=of alg] (cancel) {Kernel samples cancel the diagonal correction};
\node[box,left=of cancel] (gram) {Two weighted Gramians and a Stein identity};
\node[box,below=of gram] (norm) {\(\norm{T}\leq2\)};
\node[box,right=of norm] (limits) {Perturb \(f\), perturb \(A\), shrink domains};
\node[box,right=of limits] (main) {\(\norm{p(A)}\leq2\max_{W(A)}|p|\)};
\draw[arr] (geom)--(phi);
\draw[arr] (phi)--(cayley);
\draw[arr] (cayley)--(alg);
\draw[arr] (alg)--(cancel);
\draw[arr] (cancel)--(gram);
\draw[arr] (gram)--(norm);
\draw[arr] (norm)--(limits);
\draw[arr] (limits)--(main);
\end{tikzpicture}
\end{center}
Every arrow will be proved.  The new arrow is the passage from the
adjoint-algebra defect through exact cancellation to the Gramian
inequality.
\end{proofmap}

\section*{Deliverables and a common rubric}

The final column of the road map names one concrete artifact per week.
A ``sheet'' is a two-page collection of definitions and one checked
example; an ``audit'' annotates every change of basis or limiting step;
a ``portfolio'' collects the week's guided derivations; and a
``proof'' is a dependency-ordered argument with no unexplained
symbols.  The numerical experiment and capstone may be completed in
pairs; all other artifacts are individual.

Each deliverable is assessed on the same four questions, one point
each:
\begin{enumerate}
  \item Are all objects defined before use?
  \item Is every equality justified by a named fact or a displayed
  calculation?
  \item Are positivity, convergence, and change-of-basis claims made in
  the correct direction?
  \item Can the author state in two sentences how the artifact advances
  the main proof?
\end{enumerate}
A score of \(3/4\) is ready for the next week; a score of \(2/4\) or
below should be revised using the corresponding hints in
\cref{app:solutions}.  For the Week 12 oral proof and Week 14 full
write-up, all four criteria are required.

\mainmatter

\part{The objects in the statement}

\chapter{Matrix language: vectors, norms, adjoints, and positivity}
\label{ch:matrix-language}

\begin{learninggoals}
By the end of this week, you should be able to:
\begin{itemize}
  \item compute with the complex inner product and conjugate transpose;
  \item distinguish vector, operator, and scalar norms;
  \item translate a positive-semidefinite matrix inequality into
  inequalities of quadratic forms;
  \item explain why unitary similarity and congruence play different
  roles in the proof.
\end{itemize}
\end{learninggoals}

\begin{workshopplan}
\textbf{Opening (20 min):} compare \(Ax\), \(x^*Ax\), and eigenvalues.
\textbf{Core theory (70 min):} norms, adjoints, Hermitian matrices,
positive semidefiniteness.
\textbf{Guided proof (35 min):} congruence and square roots.
\textbf{Problem studio (45 min):} \cref{ex:norm-N,ex:psd-tests,ex:congruence}.
\end{workshopplan}

\section{Complex vectors and the adjoint}

A vector in \(\C^n\) is a column \(x=(x_1,\ldots,x_n)^T\).
Its conjugate transpose is
\[
  x^*=(\overline{x_1},\ldots,\overline{x_n}).
\]
We use the inner product and Euclidean norm
\[
  \ip{x}{y}=x^*y=\sum_{j=1}^n\overline{x_j}y_j,
  \qquad
  \norm{x}_2=\sqrt{x^*x}.
\]
This convention is conjugate-linear in the first variable and linear
in the second.  For a matrix \(A=(a_{ij})\), the adjoint is
\(A^*=\overline A^{\,T}\), characterized by
\[
  \ip{Ax}{y}=\ip{x}{A^*y}.
\]

\begin{definition}
A matrix \(U\in\Mn\) is \emph{unitary} if \(U^*U=UU^*=I\).
A matrix \(H\) is \emph{Hermitian} if \(H=H^*\).
A matrix \(A\) is \emph{normal} if \(A^*A=AA^*\).
\end{definition}

Unitary matrices preserve inner products and lengths.  They are the
complex analogue of orthogonal changes of coordinates.  The
finite-dimensional spectral theorem says that \(A\) is normal if and
only if it is unitarily diagonalizable.  Hermitian and unitary
matrices are normal; upper-triangular matrices with a nonzero
off-diagonal entry are typical nonnormal examples.

\section{The induced operator norm}

For \(A\in\Mn\), its induced Euclidean operator norm is
\[
  \norm{A}=\max_{\norm{x}_2=1}\norm{Ax}_2.
\]
The maximum exists because the unit sphere is compact and the norm is
continuous.  Three facts will be used repeatedly:
\begin{align}
  \norm{Ax}_2&\leq\norm{A}\norm{x}_2,\\
  \norm{AB}&\leq\norm{A}\norm{B},\\
  \norm{UAV}&=\norm{A}\qquad(U,V\text{ unitary}).
\end{align}
Moreover,
\[
  \norm{A}^2=\lambda_{\max}(A^*A),
\]
where \(\lambda_{\max}\) is the largest eigenvalue of a Hermitian
matrix.  Thus \(\norm{A}\) is the largest singular value, not in
general the largest modulus of an eigenvalue.

\begin{intuitionbox}
Eigenvalues measure what \(A\) does on its invariant directions.
The operator norm asks for the largest stretch in \emph{any}
direction.  For a nonnormal matrix, those can be very different.
Crouzeix's problem is fundamentally about controlling this nonnormal
stretch.
\end{intuitionbox}

\section{Positive-semidefinite order}

\begin{definition}
A Hermitian matrix \(H\) is \emph{positive semidefinite}, written
\(H\succeq0\), if
\[
  x^*Hx\geq0\qquad\text{for every }x\in\C^n.
\]
For Hermitian \(H,K\), write \(H\preceq K\) when \(K-H\succeq0\).
\end{definition}

The spectral theorem gives the equivalent conditions
\[
  H\succeq0
  \quad\Longleftrightarrow\quad
  \text{every eigenvalue of \(H\) is nonnegative}
  \quad\Longleftrightarrow\quad
  H=R^*R\text{ for some }R.
\]
Every positive-definite \(G\succ0\) has a unique positive-definite
square root \(G^{1/2}\), and \(G^{-1/2}\) is well defined.

\begin{lemma}[Congruence preserves positivity]
\label{lem:congruence}
If \(H\succeq0\), then \(S^*HS\succeq0\) for every matrix \(S\).
If \(S\) is invertible, the converse also holds.
\end{lemma}

\begin{proof}
For every \(x\),
\[
  x^*S^*HSx=(Sx)^*H(Sx)\geq0.
\]
When \(S\) is invertible, apply the same argument with \(S^{-1}\).
\end{proof}

Similarity \(S^{-1}AS\) preserves eigenvalues.  Congruence \(S^*HS\)
preserves positivity.  They coincide only for a unitary \(S\); mixing
them up is a recurring source of errors.

\section{Real and imaginary parts of a matrix}

Define
\[
  \RePart A=\frac{A+A^*}{2},
  \qquad
  \ImPart A=\frac{A-A^*}{2i}.
\]
Both are Hermitian and \(A=\RePart A+i\ImPart A\).  The expression
\(\RePart A\succeq0\) means the Hermitian part of \(A\) is positive;
it does not assert that \(A\) itself is Hermitian.

\begin{lemma}[Zero quadratic form for a positive matrix]
\label{lem:psd-zero}
If \(H\succeq0\) and \(x^*Hx=0\), then \(Hx=0\).
\end{lemma}

\begin{proof}
Write \(H=H^{1/2}H^{1/2}\).  Then
\(x^*Hx=\norm{H^{1/2}x}_2^2=0\), so \(H^{1/2}x=0\), hence \(Hx=0\).
\end{proof}

This tiny lemma is decisive in Week 12: equality in a positive Gramian
will force a vector into a kernel.

\section{Exercises}

\begin{exercise}[Warm-up: adjoints]
Verify \((AB)^*=B^*A^*\), \((A^{-1})^*=(A^*)^{-1}\), and
\(\norm{A^*}=\norm{A}\).
\end{exercise}

\begin{exercise}[Core: the sharp example]
\label{ex:norm-N}
For \(N=\begin{pmatrix}0&2\\0&0\end{pmatrix}\), compute \(N^*N\) and
show \(\norm{N}=2\).
\end{exercise}

\begin{exercise}[Core: positivity tests]
\label{ex:psd-tests}
For a Hermitian matrix
\(
H=\begin{pmatrix}a&b\\\overline b&d\end{pmatrix},
\)
prove that \(H\succeq0\) if and only if
\(a\geq0\), \(d\geq0\), and \(ad\geq|b|^2\).
\end{exercise}

\begin{exercise}[Core: congruence versus similarity]
\label{ex:congruence}
Give a \(2\times2\) example showing that similarity need not preserve
Hermitianity or positive semidefiniteness.  Then prove
\cref{lem:congruence} without referring to eigenvalues.
\end{exercise}

\begin{exercise}[Stretch: order and norm]
For Hermitian \(H\), prove
\(-cI\preceq H\preceq cI\) if and only if \(\norm{H}\leq c\).
\end{exercise}

\chapter{Spectra, diagonalization, and functions of a matrix}
\label{ch:functional-calculus}

\begin{learninggoals}
\begin{itemize}
  \item Define spectrum, resolvent, and polynomial evaluation.
  \item Explain why distinct eigenvalues imply diagonalizability.
  \item Use Lagrange interpolation to identify \(\alg(T)\).
  \item Recognize the holomorphic functional calculus, which Week 5
  develops in analytic detail.
\end{itemize}
\end{learninggoals}

\begin{workshopplan}
\textbf{Opening:} why \(p(A)\) contains more information than
\(p(\spec(A))\). \textbf{Core:} spectrum, diagonalization, generated
algebras. \textbf{Guided proof:} Lagrange interpolation.
\textbf{Studio:} calculate functions of diagonal and Jordan matrices.
\end{workshopplan}

\section{Eigenvalues, spectrum, and resolvent}

A number \(\lambda\in\C\) is an eigenvalue of \(A\) if
\(Ax=\lambda x\) for some nonzero vector \(x\).  The spectrum is
\[
  \spec(A)=\set{\lambda\in\C:\lambda I-A\text{ is not invertible}}.
\]
In finite dimensions this is exactly the set of roots of
\(\det(\lambda I-A)\).  When \(z\notin\spec(A)\), the matrix
\[
  R_A(z)=(zI-A)^{-1}
\]
is the \emph{resolvent}.

The spectral radius is
\[
  \rho(A)=\max_{\lambda\in\spec(A)}|\lambda|.
\]
Since \(Ax=\lambda x\) for a unit eigenvector,
\(\abs{\lambda}\leq\norm{A}\), so \(\rho(A)\leq\norm{A}\).
Equality need not hold.

\section{Polynomial evaluation}

If \(p(z)=a_0+a_1z+\cdots+a_dz^d\), define
\[
  p(A)=a_0I+a_1A+\cdots+a_dA^d.
\]
This definition respects sums and products:
\[
  (p+q)(A)=p(A)+q(A),\qquad (pq)(A)=p(A)q(A).
\]
The spectral mapping theorem says
\[
  \spec(p(A))=p(\spec(A)).
\]
It controls the eigenvalues of \(p(A)\), but Crouzeix's problem asks
for the generally larger quantity \(\norm{p(A)}\).

\begin{example}[A nilpotent warning]
For \(N\) from \cref{ex:norm-N}, \(\spec(N)=\set{0}\) and
\(\rho(N)=0\), yet \(\norm{N}=2\).  The eigenvalues alone give no
useful norm bound.
\end{example}

\section{Diagonalization and simple spectrum}

If \(A\) has \(n\) linearly independent eigenvectors, place them in the
columns of an invertible matrix \(S\).  Then
\[
  A=S\Lambda S^{-1},
  \qquad
  \Lambda=\diag(\lambda_1,\ldots,\lambda_n).
\]
Consequently,
\[
  p(A)=S\,\diag(p(\lambda_1),\ldots,p(\lambda_n))\,S^{-1}.
\]
Eigenvectors belonging to distinct eigenvalues are linearly
independent, so a matrix with \(n\) distinct eigenvalues is
diagonalizable.  This condition is called \emph{simple spectrum}.

\begin{pitfall}
The diagonalizing matrix \(S\) need not be unitary.  Therefore
\(\norm{S\Lambda S^{-1}}\) can be much larger than \(\norm{\Lambda}\).
Replacing a nonnormal matrix by its diagonal eigenvalue list loses the
geometry stored in \(S^*S\).
\end{pitfall}

\section{The unital algebra generated by a matrix}

\begin{definition}
The unital algebra generated by \(T\) is
\[
  \alg(T)=\set{q(T):q\in\C[z]}.
\]
\end{definition}

By Cayley--Hamilton, every high power of \(T\) is a linear combination
of \(I,T,\ldots,T^{n-1}\), so \(\alg(T)\) is finite-dimensional and
therefore norm closed.

\begin{lemma}[Interpolation description]
\label{lem:alg-simple}
Suppose
\[
  T=S\diag(\lambda_1,\ldots,\lambda_n)S^{-1}
\]
with distinct \(\lambda_j\).  Then
\[
  \alg(T)
  =\set{SDS^{-1}:D\text{ diagonal}}.
\]
Moreover,
\[
  \alg(T^*)
  =\set{(S^{-1})^*DS^*:D\text{ diagonal}}.
\]
\end{lemma}

\begin{proof}
Every polynomial in \(T\) has the first displayed form.  Conversely,
given \(D=\diag(d_1,\ldots,d_n)\), Lagrange interpolation gives a
polynomial \(q\) with \(q(\lambda_j)=d_j\), hence
\(q(T)=SDS^{-1}\).

Taking adjoints in \(T=S\Lambda S^{-1}\) gives
\[
  T^*=(S^{-1})^*\overline{\Lambda}S^*.
\]
The conjugate eigenvalues remain distinct, so the same interpolation
argument proves the second formula.
\end{proof}

The particular orientation of the second formula will drive the exact
cancellation in Week 10.

\section{First preview of the holomorphic functional calculus}

Let \(f\) be holomorphic on an open set containing \(\spec(A)\).
Because the spectrum is finite, choose disjoint small closed disks
inside that open set whose interiors contain all the distinct spectral
points, and let \(\Gamma\) be the union of their positively oriented
boundary circles.  Define
\begin{equation}
  f(A)=\frac{1}{2\pi i}\int_\Gamma
  f(z)(zI-A)^{-1}\,dz.
  \label{eq:cauchy-functional-calculus}
\end{equation}
For polynomials this agrees with the earlier definition.  If \(A\) is
diagonalizable, it gives
\[
  f(A)=S\diag(f(\lambda_1),\ldots,f(\lambda_n))S^{-1}.
\]
For a nondiagonalizable matrix, derivatives of \(f\) at eigenvalues
appear; Hermite interpolation supplies a polynomial \(q\) with
\(f(A)=q(A)\).  In particular,
\begin{equation}
  f(A)\in\alg(A)
  \quad\text{whenever \(f\) is holomorphic near \(\spec(A)\).}
  \label{eq:holomorphic-in-algebra}
\end{equation}

This is a first-pass algebraic preview.  Week 5 defines holomorphic
functions and oriented contour integrals, states Cauchy's formula, and
explains why the definition is independent of the chosen circles.
For the Week 2 exercises, it is enough to use the displayed formula or
the diagonal/Jordan formulas that follow from it.

\begin{checkpoint}
Weeks 9--10 will produce matrices \(g_{f^m}(B)\) by contour
integration.  Equation \eqref{eq:holomorphic-in-algebra} is the reason
they lie in \(\alg(B)\); taking adjoints then places the unknown
correction in \(\alg(B^*)\).
\end{checkpoint}

\section{Neumann series beyond
\texorpdfstring{\(\norm{X}<1\)}{norm X less than 1}}

The elementary geometric series says
\((I-X)^{-1}=\sum_{k\geq0}X^k\) when \(\norm{X}<1\).  In finite
dimensions it also converges whenever \(\rho(X)<1\).  One way to see
this is to use Jordan form: a Jordan block contributes a polynomial in
\(k\) times \(|\lambda|^k\), which still tends to zero and is summable
when \(|\lambda|<1\).  Therefore, if
\(\spec(T)\subset\olD\) and \(w\in\D\),
\begin{equation}
  (I-wT)^{-1}=\sum_{m=0}^{\infty}w^mT^m
  \quad\text{in operator norm}.
  \label{eq:neumann-resolvent}
\end{equation}

\section{Exercises}

\begin{exercise}[Warm-up]
Compute \(p(A)\) for
\(
A=\diag(1,i,-1)
\)
and \(p(z)=z^3-2z+1\).
\end{exercise}

\begin{exercise}[Core: interpolation]
\label{ex:lagrange}
Write the Lagrange polynomial taking prescribed values \(d_j\) at
distinct points \(\lambda_j\), and use it to prove
\cref{lem:alg-simple}.
\end{exercise}

\begin{exercise}[Core: same generated algebra]
\label{ex:same-algebra}
Let \(B\) have simple spectrum \(\beta_1,\ldots,\beta_n\), and let
\(f\) be holomorphic near \(\spec(B)\).  If
\(f(\beta_1),\ldots,f(\beta_n)\) are distinct, prove
\(\alg(f(B))=\alg(B)\).
\end{exercise}

\begin{exercise}[Core: Jordan calculus]
For \(J=\begin{pmatrix}\lambda&1\\0&\lambda\end{pmatrix}\), use
\eqref{eq:cauchy-functional-calculus} or a power series to show
\[
  f(J)=\begin{pmatrix}f(\lambda)&f'(\lambda)\\0&f(\lambda)\end{pmatrix}.
\]
\end{exercise}

\begin{exercise}[Stretch: resolvent series]
Prove \eqref{eq:neumann-resolvent} directly from Jordan normal form,
including the case of a nontrivial Jordan block.
\end{exercise}

\chapter{The numerical range: geometry attached to a matrix}
\label{ch:numerical-range}

\begin{learninggoals}
\begin{itemize}
  \item Interpret \(x^*Ax\) as a directional scalar observation of \(A\).
  \item Prove compactness, spectral containment, and unitary invariance.
  \item Understand why \(W(A)\) is convex.
  \item Read exposed faces from maximal eigenspaces of rotated Hermitian
  parts.
  \item Compute the sharp nilpotent example.
\end{itemize}
\end{learninggoals}

\begin{workshopplan}
\textbf{Opening:} sample \(x^*Ax\) numerically for a \(2\times2\)
matrix. \textbf{Core:} basic properties and the \(2\times2\) ellipse.
\textbf{Guided proof:} Toeplitz--Hausdorff by compression.
\textbf{Studio:} exact numerical ranges of elementary matrices.
\end{workshopplan}

\section{Definition and first properties}

\begin{definition}
For \(A\in\Mn\), its numerical range is
\[
  W(A)=\set{x^*Ax:x\in\C^n,\ \norm{x}_2=1}.
\]
\end{definition}

The scalar \(x^*Ax\) is the response of \(A\) observed in the same unit
direction \(x\).  The following facts are immediate but important.

\begin{proposition}
\label{prop:W-basic}
For \(A\in\Mn\):
\begin{enumerate}
  \item \(W(A)\) is compact.
  \item \(\spec(A)\subset W(A)\).
  \item \(W(U^*AU)=W(A)\) for every unitary \(U\).
  \item \(W(\alpha A+\beta I)=\alpha W(A)+\beta\) for
  \(\alpha,\beta\in\C\).
  \item \(W(A^*)=\overline{W(A)}\).
\end{enumerate}
\end{proposition}

\begin{proof}
The unit sphere is compact and \(x\mapsto x^*Ax\) is continuous,
proving (1).  A normalized eigenvector gives (2).  Substitution
\(y=Ux\) proves (3), and direct calculation proves (4)--(5).
\end{proof}

\section{The \texorpdfstring{\(2\times2\)}{2 by 2} ellipse}

\begin{theorem}[Elliptical range theorem in dimension two]
\label{thm:ellipse}
For a \(2\times2\) complex matrix \(B\), \(W(B)\) is a filled ellipse,
possibly degenerating to a disk, a line segment, or a point.  Its foci
are the two eigenvalues of \(B\).
\end{theorem}

\begin{proof}[Guided calculation]
Unitary Schur triangularization and a harmless phase change reduce to
\[
  B=\begin{pmatrix}\lambda_1&c\\0&\lambda_2\end{pmatrix},
  \qquad c\geq0.
\]
Every unit vector, up to a scalar phase that cancels from \(x^*Bx\),
has the form
\[
  x=\begin{pmatrix}\cos t\\e^{i\phi}\sin t\end{pmatrix},
  \qquad 0\leq t\leq\frac{\pi}{2}.
\]
Thus
\begin{equation}
  x^*Bx=
  \frac{\lambda_1+\lambda_2}{2}
  +\frac{\lambda_1-\lambda_2}{2}\cos(2t)
  +\frac{c}{2}e^{i\phi}\sin(2t).
  \label{eq:ellipse-param}
\end{equation}
After translating the center to \(0\) and rotating
\(\lambda_1-\lambda_2\) onto the real axis, put
\[
  s=\cos(2t),\qquad
  u=\sin(2t)\cos\phi,\qquad
  v=\sin(2t)\sin\phi.
\]
Then \((s,u,v)\) runs over the unit sphere in \(\R^3\), and the real
and imaginary coordinates in \eqref{eq:ellipse-param} are the image
of that sphere under
\[
  L(s,u,v)=\frac12\bigl(ds+cu,\;cv\bigr),
  \qquad d=|\lambda_1-\lambda_2|.
\]
The image of the unit sphere under a linear map from \(\R^3\) to
\(\R^2\) is the same as the image of the unit ball: for a point with a
preimage of norm below \(1\), add a vector in \(\ker L\) until the
preimage has norm \(1\).  The image of the ball is a filled ellipse.
Since
\[
  LL^T=\frac14
  \begin{pmatrix}d^2+c^2&0\\0&c^2\end{pmatrix},
\]
its semiaxes are
\[
  \frac12\sqrt{|\lambda_1-\lambda_2|^2+c^2}
  \quad\text{and}\quad
  \frac c2.
\]
The focal distance is \(|\lambda_1-\lambda_2|/2\), so the foci are
\(\lambda_1,\lambda_2\).
\end{proof}

\begin{blackbox}[Unitary Schur triangularization]
Every complex square matrix \(B\) admits a unitary change of basis
\(U^*BU=R\) with \(R\) upper triangular; the diagonal entries of \(R\)
are the eigenvalues of \(B\).  A standard proof chooses a unit
eigenvector, puts it first in an orthonormal basis, and proceeds by
induction on the lower-right compression.  In dimension two, a
diagonal unitary can rotate the single off-diagonal entry of \(R\) to
a nonnegative real number.  This is exactly the reduction used above.
\end{blackbox}

\section{Toeplitz--Hausdorff convexity}

\begin{theorem}[Toeplitz--Hausdorff]
\label{thm:toeplitz-hausdorff}
The numerical range of every complex matrix is convex.
\end{theorem}

\begin{proof}
Take \(z_1,z_2\in W(A)\), realized by unit vectors \(x_1,x_2\).
Let \(\mathcal M=\Span\{x_1,x_2\}\) and let
\(B=P_{\mathcal M}A|_{\mathcal M}\) be the compression of \(A\) to
this subspace.  If \(\dim\mathcal M=1\), then \(x_1\) and \(x_2\)
differ only by a scalar phase, so \(z_1=z_2\) and there is nothing to
prove.  Otherwise \(\dim\mathcal M=2\).  For every unit
\(x\in\mathcal M\),
\[
  x^*Bx=x^*P_{\mathcal M}Ax=x^*Ax,
\]
so \(W(B)\subset W(A)\), while \(z_1,z_2\in W(B)\).
By \cref{thm:ellipse}, \(W(B)\) is convex and therefore contains the
line segment joining \(z_1\) and \(z_2\).  The same segment lies in
\(W(A)\).
\end{proof}

\section{Support lines and Hermitian parts}

For \(\theta\in\R\), put
\[
  H_\theta=\RePart(e^{-i\theta}A)
  =\frac{e^{-i\theta}A+e^{i\theta}A^*}{2}.
\]
Then
\begin{equation}
  \max_{z\in W(A)}\RePart(e^{-i\theta}z)
  =\lambda_{\max}(H_\theta).
  \label{eq:W-support-function}
\end{equation}
Indeed, the left side is
\[
  \max_{\norm{x}=1}x^*H_\theta x,
\]
which is the Rayleigh-quotient characterization of the largest
Hermitian eigenvalue.  Formula \eqref{eq:W-support-function} both
explains the geometry and gives a practical way to plot \(W(A)\).

\begin{proposition}[Exposed faces from the support pencil]
\label{prop:exposed-face-pencil}
Fix \(\theta\in\R\), and write
\[
  h_\theta=\lambda_{\max}(H_\theta),
  \qquad
  E_\theta=\ker(h_\theta I-H_\theta).
\]
Let \(P_\theta\) be the orthogonal projection onto \(E_\theta\), and
define the Hermitian tangential compression
\[
  K_\theta
  =
  P_\theta\ImPart(e^{-i\theta}A)|_{E_\theta}.
\]
The exposed face of \(W(A)\) with outward normal \(e^{i\theta}\) is
\begin{equation}
  \set{z\in W(A):\RePart(e^{-i\theta}z)=h_\theta}
  =
  e^{i\theta}
  \left(
    h_\theta+
    i[\lambda_{\min}(K_\theta),\lambda_{\max}(K_\theta)]
  \right).
  \label{eq:exposed-face-pencil}
\end{equation}
\end{proposition}

\begin{proof}
The support slack \(h_\theta I-H_\theta\) is positive semidefinite.  If
a unit vector \(x\) generates a point on the exposed face, then
\[
  x^*(h_\theta I-H_\theta)x=0.
\]
By \cref{lem:psd-zero}, this is equivalent to \(x\in E_\theta\).
Conversely, every unit vector in \(E_\theta\) generates a point on the
supporting line.  For such a vector,
\[
  e^{-i\theta}x^*Ax
  =
  h_\theta+
  i\,x^*\ImPart(e^{-i\theta}A)x.
\]
The Rayleigh quotients of \(K_\theta\) fill the interval between its
smallest and largest eigenvalues, which proves
\eqref{eq:exposed-face-pencil}.
\end{proof}

\begin{intuitionbox}[The planar boundary hides a spectral bundle]
A simple top eigenvalue of \(H_\theta\) produces one exposed point.
A non-scalar \(K_\theta\) produces a nontrivial flat edge.  Thus the
nullspace of the support slack records all vector states that generate
the supported face; the planar outline shows only their scalar
expectations.
\end{intuitionbox}

\section{The example that forces the constant
\texorpdfstring{\(2\)}{2}}

Let
\[
  N=\begin{pmatrix}0&2\\0&0\end{pmatrix}.
\]
For
\(
x=(\cos t,e^{i\phi}\sin t)^T
\),
\[
  x^*Nx=e^{i\phi}\sin(2t).
\]
As \(t\) and \(\phi\) vary, this fills \(\olD\).  Hence
\[
  W(N)=\olD.
\]
From Week 1, \(\norm{N}=2\).  Taking \(p(z)=z\) gives
\[
  \norm{p(N)}=2
  =2\max_{z\in W(N)}|p(z)|.
\]
No universal constant smaller than \(2\) can work.

\section{Continuity of the numerical range}

\begin{lemma}[Lipschitz inclusion]
\label{lem:W-Lipschitz}
For \(A,B\in\Mn\),
\[
  W(B)\subset W(A)+\set{z\in\C:|z|\leq\norm{B-A}}.
\]
\end{lemma}

\begin{proof}
For every unit vector \(x\),
\[
  |x^*(B-A)x|
  \leq\norm{(B-A)x}_2\norm{x}_2
  \leq\norm{B-A}.
\]
\end{proof}

This estimate will keep numerical ranges inside a fixed outer domain
while matrices with simple spectrum approach an arbitrary matrix.

\section{Exercises}

\begin{exercise}[Warm-up]
Compute \(W(A)\) when \(A\) is Hermitian.  Then compute it when
\(A=\lambda I\).
\end{exercise}

\begin{exercise}[Core: nilpotent disk]
\label{ex:nilpotent-disk}
Fill in every step showing \(W(N)=\olD\) for the matrix above.
\end{exercise}

\begin{exercise}[Core: support function]
Prove \eqref{eq:W-support-function}.  Explain how a largest-eigenvalue
routine can trace boundary points of \(W(A)\).
\end{exercise}

\begin{exercise}[Core: compression]
If \(\mathcal M\subset\C^n\) and \(B=P_{\mathcal M}A|_{\mathcal M}\),
prove \(W(B)\subset W(A)\).  Why is this the correct direction for the
Toeplitz--Hausdorff proof?
\end{exercise}

\begin{exercise}[Stretch: ellipse details]
Verify the computation of \(LL^T\) after
\eqref{eq:ellipse-param} and recover the two semiaxis lengths.
Alternatively, obtain the same ellipse from its support function.
\end{exercise}

\chapter{Spectral sets and Crouzeix's problem}
\label{ch:spectral-sets}

\begin{learninggoals}
\begin{itemize}
  \item State the notion of a \(C\)-spectral set.
  \item Explain why numerical ranges are natural domains for matrix
  functions.
  \item Separate the normal case from the nonnormal difficulty.
  \item Normalize the main inequality to a unit-bounded function.
\end{itemize}
\end{learninggoals}

\begin{workshopplan}
\textbf{Opening:} compare the spectrum and numerical range of the
sharp nilpotent example. \textbf{Core:} spectral-set language and the
constant question. \textbf{Guided proof:} normalize a polynomial to
unit supremum. \textbf{Studio:} write the one-page problem brief,
including the normal and nilpotent benchmark cases.
\end{workshopplan}

\section{What a spectral set controls}

\begin{definition}
Let \(K\subset\C\) be compact and contain \(\spec(A)\).  It is a
\(C\)-spectral set for \(A\) if
\[
  \norm{r(A)}
  \leq C\max_{z\in K}|r(z)|
\]
for every rational function \(r\) with no pole on \(K\).
\end{definition}

Here rational matrix evaluation is concrete.  Write \(r=p/q\) with
polynomials \(p,q\).  Having no pole on \(K\supset\spec(A)\) means
\(q(\lambda)\neq0\) for every \(\lambda\in\spec(A)\); spectral mapping
then makes \(q(A)\) invertible, and
\[
  r(A)=p(A)q(A)^{-1}.
\]
The value is independent of the chosen fraction and agrees with the
holomorphic functional calculus.

The scalar maximum on the right is often accessible geometrically or
numerically.  The operator norm on the left measures the actual
amplification of the matrix function.  A spectral-set theorem transfers
a scalar approximation guarantee to a matrix approximation guarantee.

\begin{example}[Normal matrices]
If \(A\) is normal, it is unitarily diagonalizable.  Therefore
\[
  \norm{p(A)}
  =\max_{\lambda\in\spec(A)}|p(\lambda)|
  \leq\max_{z\in W(A)}|p(z)|.
\]
Thus \(W(A)\) gives constant \(1\) for normal matrices.  The universal
constant is forced above \(1\) only by nonnormal matrices.
\end{example}

\section{The main theorem}

\begin{theorem}[Main theorem of the source manuscript]
\label{thm:main}
For every \(A\in\Mn\) and every polynomial \(p\in\C[z]\),
\begin{equation}
  \norm{p(A)}
  \leq2\max_{z\in W(A)}|p(z)|.
  \label{eq:main}
\end{equation}
Consequently, \(W(A)\) is a \(2\)-spectral set for \(A\).
\end{theorem}

The factor \(2\) cannot be lowered, by the nilpotent example in Week 3.
Historically, Crouzeix formulated the conjecture and proved foundational
bounds \cite{Crouzeix2004,Crouzeix2007}; Crouzeix and Palencia obtained
the universal constant \(1+\sqrt2\)
\cite{CrouzeixPalencia2017}.  The source manuscript's new mechanism
is claimed to reach \(2\).

\section{Normalization}

Fix a bounded domain \(\Omega\) containing \(W(A)\), and set
\[
  m_\Omega=\max_{z\in\overline\Omega}|p(z)|.
\]
If \(m_\Omega>0\), let \(f=p/m_\Omega\).  Then
\[
  \max_{\overline\Omega}|f|\leq1,
  \qquad
  p(A)=m_\Omega f(A).
\]
Thus the core problem becomes:
\[
  \boxed{\quad |f|\leq1\text{ on the domain}
  \quad\Longrightarrow\quad \norm{f(A)}\leq2.\quad}
\]
This unit-ball normalization is what makes the Cayley transform in
Week 9 have positive real part.

\section{Why the numerical range is the right set}

The spectrum is too small for nonnormal matrices: the nilpotent \(N\)
has spectrum \(\{0\}\) but norm \(2\).  A large disk based only on
\(\norm{A}\) is often too crude.  The numerical range sits between:
\[
  \spec(A)\subset W(A)\subset\set{z:|z|\leq\norm{A}},
\]
is compact and convex, changes continuously with \(A\), and records
directional information \(x^*Ax\) that the spectrum forgets.

\section{Polynomial, holomorphic, and rational versions}

The polynomial estimate is the core.  Once it is known, approximation
will give:
\begin{itemize}
  \item the same bound for functions holomorphic near \(W(A)\);
  \item the rational-function definition of a spectral set;
  \item a finite-dimensional-compression version for bounded operators
  on a Hilbert space.
\end{itemize}
These passages are deferred to Weeks 14--15 so that approximation does
not obscure the new finite-dimensional argument.

\begin{proofmap}[Four layers of the proof]
\begin{enumerate}
  \item \textbf{Geometry:} put the numerical range in a smooth convex
  outer domain.
  \item \textbf{Positivity:} use the double-layer measure and Cayley
  transforms to build a positive-real matrix function.
  \item \textbf{Algebraic de-symmetrization:} cancel the adjoint-algebra
  defect and prove the norm bound \(2\).
  \item \textbf{Approximation:} remove simple spectrum, algebra
  generation, and the outer domain.
\end{enumerate}
\end{proofmap}

\section{Exercises}

\begin{exercise}[Warm-up]
Show that if \(K\) is a \(C\)-spectral set for \(A\), then
\(\alpha K+\beta\) is a \(C\)-spectral set for \(\alpha A+\beta I\)
when \(\alpha\neq0\).
\end{exercise}

\begin{exercise}[Core: normal case]
Prove carefully that a normal matrix satisfies \eqref{eq:main} with
constant \(1\).
\end{exercise}

\begin{exercise}[Core: normalization]
Explain separately what happens when \(m_\Omega=0\), and justify why
normalizing by \(m_\Omega\) loses no information.
\end{exercise}

\begin{exercise}[Stretch: approximation error]
Assume \cref{thm:main}.  If \(f\) is holomorphic near \(W(A)\) and \(s\)
is a polynomial, derive
\[
  \norm{f(A)-s(A)}
  \leq2\max_{z\in W(A)}|f(z)-s(z)|.
\]
List the additional fact needed to define \(f(A)\).
\end{exercise}

\part{Analytic and positive-function machinery}

\chapter{Complex analysis and convex boundary geometry}
\label{ch:complex-convex}

\begin{learninggoals}
\begin{itemize}
  \item Read the contour integrals used in the functional calculus.
  \item Use Cauchy's formula and the maximum-modulus principle.
  \item Track boundary orientation, tangent, and outward normal.
  \item State the exact approximation theorem needed at the end.
\end{itemize}
\end{learninggoals}

\begin{workshopplan}
\textbf{Opening:} orient a circle and locate its outward normal.
\textbf{Core:} holomorphic functions, Cauchy formula, matrix-valued
integration. \textbf{Geometry lab:} supporting lines and
\(d\zeta=i\nu\,ds\). \textbf{Closing:} Runge approximation as a
clearly marked classical input.
\end{workshopplan}

\section{Holomorphic functions and contour integrals}

A function \(f:\Omega\to\C\), on an open set \(\Omega\), is
\emph{holomorphic} if its complex derivative exists at every point.
Holomorphic functions are far more rigid than differentiable real
functions: they have local power series, obey Cauchy's integral
formula, and cannot attain a strict interior maximum of their modulus
unless they are constant.

If \(\Gamma\) is a positively oriented simple closed \(C^1\) curve and
\(f\) is holomorphic on a neighborhood of the region it encloses, then
\begin{equation}
  f(z)=\frac{1}{2\pi i}\int_\Gamma
  \frac{f(\zeta)}{\zeta-z}\,d\zeta
  \qquad(z\text{ inside }\Gamma).
  \label{eq:scalar-cauchy}
\end{equation}
Applying this entry by entry gives the matrix Cauchy formula.  In
particular, if \(\spec(B)\) lies inside \(\Gamma\),
\begin{equation}
  \frac{1}{2\pi i}\int_\Gamma
  (\zeta I-B)^{-1}\,d\zeta=I.
  \label{eq:matrix-resolvent-mass}
\end{equation}

\begin{blackbox}[Cauchy's theorem and formula]
The standard proof first establishes that integrals of holomorphic
functions around triangles vanish, deforms contours through a
triangulation, and applies the result to
\((f(\zeta)-f(z))/(\zeta-z)\).  In this workshop, the scalar formula
\eqref{eq:scalar-cauchy} is a classical input; every matrix use is a
finite collection of scalar uses.
\end{blackbox}

\section{Maximum modulus}

\begin{theorem}[Maximum-modulus principle]
If \(f\) is holomorphic on a neighborhood of the closure of a bounded
domain \(\Omega\), then
\[
  \max_{z\in\overline\Omega}|f(z)|
  =\max_{\zeta\in\partial\Omega}|f(\zeta)|.
\]
\end{theorem}

For the proof, the important consequence is that a boundary bound
\(|f|\leq1\) controls the entire domain.  We often state the bound on
\(\overline\Omega\) to avoid invoking the principle each time.

\section{Smooth convex domains and normals}

Let \(\Omega\subset\C\cong\R^2\) be bounded and convex, with
continuously differentiable boundary.  At
\(\zeta\in\partial\Omega\), let \(\nu(\zeta)\in\C\) be the outward
unit normal and let \(ds\) denote arclength.  When the boundary is
traversed counterclockwise, the unit tangent is \(i\nu\), so
\begin{equation}
  d\zeta=i\nu(\zeta)\,ds,
  \qquad
  d\overline\zeta=-i\overline{\nu(\zeta)}\,ds.
  \label{eq:orientation}
\end{equation}

Convexity says the tangent line supports the whole domain.  In complex
notation,
\begin{equation}
  \RePart\bigl(\overline{\nu(\zeta)}(\zeta-z)\bigr)\geq0
  \qquad
  (z\in\overline\Omega,\ \zeta\in\partial\Omega).
  \label{eq:supporting-line}
\end{equation}
The expression is simply the outward normal component of the vector
from \(z\) to \(\zeta\).

\begin{example}[The disk]
For \(\Omega=\set{z:|z|<R}\), write
\(\zeta=Re^{it}\).  Then \(\nu=e^{it}\), \(ds=R\,dt\), and
\[
  i\nu\,ds=iRe^{it}\,dt=d\zeta.
\]
The support inequality becomes
\(\RePart(R-e^{-it}z)\geq0\), which follows from \(|z|\leq R\).
\end{example}

\begin{pitfall}[Orientation signs]
The identities for \(d\zeta\) and \(d\overline\zeta\) have opposite
signs.  The second sign is needed when the companion integral is
adjointed in Week 8.
\end{pitfall}

\section{Uniform convergence and matrix evaluation}

Suppose \(f_j\to f\) uniformly on a contour \(\Gamma\) enclosing
\(\spec(A)\).  From the Cauchy functional calculus,
\[
  \norm{f_j(A)-f(A)}
  \leq \frac{\operatorname{length}(\Gamma)}{2\pi}
  \max_{\zeta\in\Gamma}|f_j(\zeta)-f(\zeta)|
  \max_{\zeta\in\Gamma}\norm{(\zeta I-A)^{-1}}.
\]
Thus \(f_j(A)\to f(A)\).  This elementary estimate is the bridge from
scalar approximation to matrix approximation.

\begin{blackbox}[Runge approximation in the needed form]
Let \(K\subset\C\) be compact and suppose \(\C\setminus K\) is
connected.  If \(f\) is holomorphic on a neighborhood of \(K\), then
there are polynomials \(p_j\) converging uniformly to \(f\) on \(K\).

We use this only when \(K\) is a compact convex neighborhood of a
numerical range.  Convexity makes the complement connected.  A standard
proof first approximates contour integrals by rational functions and
then moves all poles to infinity when the complement is connected.
\end{blackbox}

\section{Distance, Hausdorff convergence, and outer domains}

For a point \(z\) and set \(K\),
\[
  \dist(z,K)=\inf_{w\in K}|z-w|.
\]
For nonempty compact sets \(K,L\), their Hausdorff distance is
\[
  d_{\mathrm H}(K,L)=
  \max\left\{\sup_{z\in K}\dist(z,L),
             \sup_{w\in L}\dist(w,K)\right\}.
\]
The \(\delta\)-parallel body
\[
  K+\delta\olD=\set{k+\delta u:k\in K,\ |u|\leq1}
\]
lies at Hausdorff distance \(\delta\) from \(K\).  When \(K\) is
convex, so is this parallel body.  Week 13 will prove the regularity
needed to use its interior as a double-layer domain.

\section{Exercises}

\begin{exercise}[Warm-up: orientation]
Verify \eqref{eq:orientation} for an arclength parametrization
\(\gamma(s)\) of a counterclockwise \(C^1\) curve.
\end{exercise}

\begin{exercise}[Core: matrix Cauchy mass]
\label{ex:matrix-cauchy-mass}
Prove \eqref{eq:matrix-resolvent-mass} first for a diagonalizable
matrix and then explain why the holomorphic functional calculus gives
it for every matrix.
\end{exercise}

\begin{exercise}[Core: supporting line]
Translate \eqref{eq:supporting-line} into a dot-product inequality in
\(\R^2\), and prove it for an arbitrary convex set with a differentiable
boundary.
\end{exercise}

\begin{exercise}[Core: approximation transfer]
Derive the displayed estimate for
\(\norm{f_j(A)-f(A)}\) and identify every constant.
\end{exercise}

\begin{exercise}[Stretch: Hausdorff maxima]
Let \(K_j,K\) be nonempty compact sets contained in one compact set,
with \(d_{\mathrm H}(K_j,K)\to0\).  Prove that for every continuous
\(\varphi\),
\[
  \max_{K_j}\varphi\longrightarrow\max_K\varphi.
\]
\end{exercise}

\chapter{Positive matrices, measures, and maps}
\label{ch:positive-maps}

\begin{learninggoals}
\begin{itemize}
  \item Define a positive matrix-valued measure and integrate against it.
  \item Recognize a positive and a unital positive map.
  \item Move positivity through integration.
  \item Understand the scalar Cayley transform as a positivity
  generator.
\end{itemize}
\end{learninggoals}

\begin{workshopplan}
\textbf{Opening:} turn a two-point scalar measure into a diagonal
matrix measure. \textbf{Core:} positivity of measures and maps.
\textbf{Guided proof:} matrix-measure Cauchy--Schwarz and
contractivity. \textbf{Studio:} audit exactly where unitality and
positivity are used.
\end{workshopplan}

\section{From scalar measures to matrix-valued measures}

A finite positive scalar measure \(\mu\) assigns a nonnegative mass
\(\mu(E)\) to Borel sets and permits integration of continuous
functions.  A finite \(\Mn\)-valued measure \(M\) assigns a matrix
\(M(E)\) to each Borel set, countably additively entry by entry.
It is \emph{positive} if
\[
  M(E)\succeq0\qquad\text{for every Borel set }E.
\]
Equivalently, \(u^*M(\,\cdot\,)u\) is a positive scalar measure for
every \(u\in\C^n\).

For a continuous scalar function \(\varphi\), the matrix integral is
defined entry by entry.  Positivity gives
\begin{equation}
  \varphi\geq0
  \quad\Longrightarrow\quad
  \int\varphi\,dM\succeq0,
  \label{eq:positive-integral}
\end{equation}
because its quadratic form on \(u\) is the scalar integral against
\(u^*M(\,\cdot\,)u\).

\begin{example}[A discrete positive matrix measure]
Choose points \(\zeta_1,\ldots,\zeta_m\) and matrices
\(M_k\succeq0\).  Then
\[
  M=\sum_{k=1}^m M_k\delta_{\zeta_k}
\]
is positive, and
\(
\int\varphi\,dM=\sum_k\varphi(\zeta_k)M_k.
\)
\end{example}

\section{Positive linear maps}

For a compact space \(X\), \(C(X)\) denotes the complex-valued
continuous functions on \(X\), with pointwise addition and
multiplication and the supremum norm
\(\norm{\varphi}_\infty=\max_{x\in X}|\varphi(x)|\).

\begin{definition}
A complex-linear map
\(\Phi:C(X)\to\Mn\), where \(X\) is compact, is \emph{positive} if
\[
  \varphi\geq0\text{ pointwise}
  \quad\Longrightarrow\quad
  \Phi(\varphi)\succeq0.
\]
It is \emph{unital} if \(\Phi(1)=I\).
\end{definition}

If \(M\) is a positive matrix measure of total mass \(I\), then
\[
  \Phi(\varphi)=\int_X\varphi\,dM
\]
is unital and positive.  Positive maps are automatically
\(*\)-preserving:
\[
  \Phi(\overline\varphi)=\Phi(\varphi)^*.
\]
Indeed, a real-valued continuous function is a difference of two
nonnegative functions, so its image is Hermitian; split a complex
function into real and imaginary parts.

\begin{lemma}[Contractivity of an integral positive map]
\label{lem:positive-map-bounded}
Let \(M\) be a positive \(\Mn\)-valued measure on \(X\) with
\(M(X)=I\), and define
\(\Phi(\varphi)=\int_X\varphi\,dM\).  Then
\[
  \norm{\Phi(\varphi)}\leq\norm{\varphi}_\infty.
\]
\end{lemma}

\begin{proof}
For continuous vector-valued functions \(F,G:X\to\C^n\), set
\[
  [F,G]_M=\int_X F(\zeta)^*\,dM(\zeta)\,G(\zeta).
\]
Positivity of \(M\) makes this a positive semidefinite sesquilinear
form, so the usual quadratic-polynomial proof gives Cauchy--Schwarz:
\[
  |[F,G]_M|^2\leq[F,F]_M[G,G]_M.
\]
For unit vectors \(x,y\), choose
\(
F(\zeta)=\overline{\varphi(\zeta)}x
\)
and \(G(\zeta)=y\).  Then
\[
  |x^*\Phi(\varphi)y|
  \leq
  \left(\int_X|\varphi|^2\,x^*dM\,x\right)^{1/2}
  \left(\int_Xy^*dM\,y\right)^{1/2}
  \leq\norm{\varphi}_\infty.
\]
Taking the supremum over unit \(x,y\) proves the claim.
\end{proof}

\begin{pitfall}
A positive map is generally not multiplicative:
\(\Phi(f^2)\) need not equal \(\Phi(f)^2\).  The proof keeps
\(\Phi(f^m)\) as separate coefficients; it never makes this invalid
replacement.
\end{pitfall}

\section{The scalar Cayley transform}

For \(|a|\leq1\) and \(w\in\D\), define
\[
  c_w(a)=\frac{1+wa}{1-wa}.
\]
The denominator is nonzero and
\begin{equation}
  \RePart c_w(a)
  =\frac{1-|wa|^2}{|1-wa|^2}\geq0.
  \label{eq:cayley-real}
\end{equation}
Also,
\begin{equation}
  c_w(a)=1+2\sum_{m=1}^{\infty}w^ma^m.
  \label{eq:cayley-geometric}
\end{equation}
Thus the Cayley transform packages every power of \(a\) into one
analytic function with positive real part.

\begin{intuitionbox}
The unit disk is converted into the right half-plane.  Disk bounds are
hard nonlinear inequalities; positive real part is a linear
matrix-order condition after applying a positive map.  Retaining the
whole parameter \(w\) preserves all power coefficients.
\end{intuitionbox}

\section{Exercises}

\begin{exercise}[Warm-up]
Prove \eqref{eq:cayley-real} by multiplying numerator and denominator
by \(1-\overline{wa}\).
\end{exercise}

\begin{exercise}[Core: positive integration]
Prove \eqref{eq:positive-integral} by testing on an arbitrary vector.
\end{exercise}

\begin{exercise}[Core: star preservation]
Give the complete proof that a positive complex-linear map
\(\Phi:C(X)\to\Mn\) is \(*\)-preserving.
\end{exercise}

\begin{exercise}[Core: the factor two]
Derive \eqref{eq:cayley-geometric}.  Explain why the coefficient \(2\)
will compensate for the \(1/2\) in the definition of the double-layer
map.
\end{exercise}

\begin{exercise}[Stretch]
Prove the Cauchy--Schwarz inequality for the form
\([F,G]_M\) used in \cref{lem:positive-map-bounded}, including the
case in which \([G,G]_M=0\).
\end{exercise}

\chapter{Positive-real functions and the Herglotz kernel}
\label{ch:herglotz}

\begin{learninggoals}
\begin{itemize}
  \item Define a matrix-valued positive kernel.
  \item State scalar and matrix Herglotz representations.
  \item Prove that a positive-real analytic matrix function produces a
  positive kernel.
  \item Treat finite kernel sampling as an inequality-generating tool.
\end{itemize}
\end{learninggoals}

\begin{workshopplan}
\textbf{Opening:} scalar Cayley kernels.
\textbf{Core:} positive kernels and scalar Herglotz.
\textbf{Guided proof:} polarize scalar measures into a matrix measure.
\textbf{Studio:} two-point and multi-point sampling.
\end{workshopplan}

\section{Positive kernels}

\begin{definition}
An \(\Mn\)-valued kernel \(\mathcal K\) on \(X\times X\) is
\emph{positive} if, for every finite choice of points
\(x_0,\ldots,x_m\in X\) and vectors
\(\xi_0,\ldots,\xi_m\in\C^n\),
\begin{equation}
  \sum_{i,j=0}^m
  \xi_i^*\mathcal K(x_i,x_j)\xi_j\geq0.
  \label{eq:positive-kernel}
\end{equation}
\end{definition}

For scalar kernels this says every finite Gram matrix
\([K(x_i,x_j)]\) is positive semidefinite.  For matrix-valued kernels,
the vectors \(\xi_i\) select both a point and a direction.

\begin{example}[A factored kernel]
If \(V:X\to\mathbb M_{r,n}(\C)\), then
\[
  \mathcal K(x,y)=V(x)^*V(y)
\]
is positive, since the left side of \eqref{eq:positive-kernel} is
\(
\norm{\sum_jV(x_j)\xi_j}_2^2.
\)
Integral mixtures of such kernels are positive as well.
\end{example}

\section{The scalar Herglotz representation}

\begin{blackbox}[Scalar Herglotz theorem]
If \(f:\D\to\C\) is holomorphic and \(\RePart f(w)\geq0\), then there
are a finite positive measure \(\mu\) on \(\Torus\) and a real number
\(c\) such that
\begin{equation}
  f(w)=ic+\int_{\Torus}\frac{\zeta+w}{\zeta-w}\,d\mu(\zeta).
  \label{eq:scalar-herglotz}
\end{equation}
They are unique.

Proof blueprint: the positive harmonic function \(\RePart f\) is the
Poisson integral of a boundary measure obtained as a weak limit of its
values on circles; a harmonic conjugate recovers \(f\), up to the
imaginary constant \(ic\).  Uniqueness follows from uniqueness of
Fourier coefficients.
\end{blackbox}

\section{Assembling the matrix representation}

\begin{lemma}[Matrix Herglotz representation]
\label{lem:matrix-herglotz-representation}
Let \(F:\D\to\Mn\) be analytic with
\(\RePart F(w)\succeq0\).  Then there are a positive
\(\Mn\)-valued measure \(M\) on \(\Torus\) and a Hermitian matrix
\(J_F\) such that
\begin{equation}
  F(w)=iJ_F+
  \int_{\Torus}\frac{\zeta+w}{\zeta-w}\,dM(\zeta).
  \label{eq:matrix-herglotz}
\end{equation}
\end{lemma}

\begin{proof}
For each \(u\in\C^n\), the scalar function \(u^*F(w)u\) has
nonnegative real part.  By the scalar Herglotz theorem, it has a
positive measure \(\mu_u\) and a real constant \(c_u\).
Uniqueness in the scalar theorem gives the measure identities
\[
  \mu_{\alpha u}=|\alpha|^2\mu_u,
  \qquad
  \mu_{u+v}+\mu_{u-v}=2\mu_u+2\mu_v.
\]
Thus \(u\mapsto\mu_u(E)\) is a nonnegative quadratic form for each
Borel set \(E\).  Polarization defines
\begin{equation}
  x^*M(E)y=
  \frac14\sum_{k=0}^3 i^{-k}\,
  \mu_{x+i^ky}(E).
  \label{eq:measure-polarization}
\end{equation}
Then \(u^*M(E)u=\mu_u(E)\geq0\), so \(M\) is positive.  The imaginary
constants polarize to
\[
  J_F=\frac{F(0)-F(0)^*}{2i},
\]
which is Hermitian.  Since all quadratic forms of the two sides agree,
polarization assembles the scalar formulas into
\eqref{eq:matrix-herglotz}.
\end{proof}

\section{The Herglotz kernel}

\begin{lemma}[Matrix Herglotz kernel]
\label{lem:herglotz-kernel}
If \(F:\D\to\Mn\) is analytic and
\(\RePart F(w)\succeq0\), then
\begin{equation}
  \mathcal L_F(w,z)
  =\frac{F(w)+F(z)^*}{1-w\overline z}
  \label{eq:herglotz-kernel}
\end{equation}
is a positive kernel on \(\D\).
\end{lemma}

\begin{proof}
For \(|\zeta|=1\), direct algebra gives
\begin{equation}
 \frac{1}{1-w\overline z}
 \left(
   \frac{\zeta+w}{\zeta-w}
   +\overline{\frac{\zeta+z}{\zeta-z}}
 \right)
 =
 \frac{2}{(1-w\overline\zeta)(1-\overline z\zeta)}.
 \label{eq:herglotz-factor}
\end{equation}
The term \(iJ_F\) cancels when \(F(w)+F(z)^*\) is formed.
Substitution of \eqref{eq:matrix-herglotz} gives a factored integral.
For finite samples \(w_i,\xi_i\), the kernel quadratic form equals
\[
  2\int_{\Torus}
  \left(\sum_i\frac{\xi_i}{1-\overline{w_i}\zeta}\right)^*
  dM(\zeta)
  \left(\sum_j\frac{\xi_j}{1-\overline{w_j}\zeta}\right)
  \geq0.
\]
This is precisely \eqref{eq:positive-kernel}.
\end{proof}

\begin{intuitionbox}[A positive function becomes many inequalities]
The pointwise condition \(\RePart F(w)\succeq0\) is already useful.
The Herglotz kernel is stronger operationally: every deliberate finite
choice of points and vectors produces a scalar inequality.  The new
proof succeeds by choosing samples that annihilate what is unknown and
retain what is useful.
\end{intuitionbox}

\section{Sampling as test design}

The definition permits repeated points, zero vectors, and points
chosen from spectral data.  One should think of
\(\{(w_i,\xi_i)\}\) as a designed experiment:
\begin{enumerate}
  \item split \(F=F_{\mathrm{known}}+F_{\mathrm{nuisance}}\);
  \item choose samples that make the nuisance quadratic form zero;
  \item read a positive inequality from the known part.
\end{enumerate}
Nothing says the kernel matrix of the known part alone is positive for
all samples.  Positivity is asserted only for the original kernel, and
the cancellation works only for the special test data.

\section{Exercises}

\begin{exercise}[Warm-up: factorization]
Prove that every kernel \(V(x)^*V(y)\) is positive.
\end{exercise}

\begin{exercise}[Core: scalar identity]
\label{ex:herglotz-factor}
Verify \eqref{eq:herglotz-factor}, keeping track of
\(\overline\zeta=\zeta^{-1}\).
\end{exercise}

\begin{exercise}[Core: polarization]
Check \eqref{eq:measure-polarization} and show that its diagonal value
is \(\mu_u(E)\).
\end{exercise}

\begin{exercise}[Core: two samples]
Write the kernel inequality explicitly for points \(0,w\) and vectors
\(u,v\).  Identify the \(2\times2\) block matrix whose quadratic form
appears.
\end{exercise}

\begin{exercise}[Stretch: direct Cayley atom]
For \(F(w)=(1+wA)(1-wA)^{-1}\) with a scalar \(|A|\leq1\), factor its
Herglotz kernel directly and compare with
\eqref{eq:herglotz-factor}.
\end{exercise}

\chapter{The classical double-layer positive map}
\label{ch:double-layer}

\begin{learninggoals}
\begin{itemize}
  \item Construct the double-layer boundary measure.
  \item Prove its positivity from convex support lines.
  \item Recognize positivity as an exact support-slack test for
  numerical-range containment.
  \item Prove its total mass is \(2I\).
  \item Derive the companion identity
  \(2\Phi(h)=h(B)+g_h(B)^*\).
  \item Interpret the normalized measure as a boundary POVM and a
  symmetrized normal dilation.
\end{itemize}
\end{learninggoals}

\begin{workshopplan}
\textbf{Opening:} support lines become positive quadratic forms.
\textbf{Core derivation:} resolvent density, positivity, and mass.
\textbf{Contour lab:} companion transform with an interior contour.
\textbf{Closing comparison:} what the classical symmetrized identity
does and does not control.
\end{workshopplan}

\section{Setup}

Let \(B\in\Mn\), and let \(\Omega\subset\C\) be a bounded convex
domain with \(C^1\) boundary such that
\[
  W(B)\subset\Omega.
\]
For \(\zeta\in\partial\Omega\), define
\[
  Z_\zeta=\zeta I-B,
  \qquad
  \mathcal R_\zeta=Z_\zeta^{-1}.
\]
The resolvent exists because
\(\spec(B)\subset W(B)\subset\Omega\).

\section{The boundary density is positive}

Define the \(\Mn\)-valued boundary measure
\begin{equation}
  d\mu_B(\zeta)=\frac{1}{2\pi}
  \left(
    \nu(\zeta)\mathcal R_\zeta+
    \overline{\nu(\zeta)}\mathcal R_\zeta^*
  \right)ds.
  \label{eq:double-layer-measure}
\end{equation}
We show the matrix density in parentheses is positive.

For a unit vector \(x\), the supporting-line inequality gives
\[
  2\RePart\left(
  \overline{\nu(\zeta)}(\zeta-x^*Bx)\right)\geq0.
\]
Equivalently,
\begin{equation}
  \nu(\zeta)Z_\zeta^*
  +\overline{\nu(\zeta)}Z_\zeta\succeq0.
  \label{eq:support-matrix}
\end{equation}
Congruence on the left by \(\mathcal R_\zeta^*\) and on the right by
\(\mathcal R_\zeta\) turns the left side into
\[
  \nu(\zeta)\mathcal R_\zeta+
  \overline{\nu(\zeta)}\mathcal R_\zeta^*.
\]
Thus \(d\mu_B\) is a positive matrix measure.

\begin{intuitionbox}
Convexity says every point \(x^*Bx\) lies behind the tangent line.
The resolvent congruence lifts that scalar geometric fact into a
positive matrix density.  This is where numerical-range geometry
enters the analytic machinery.
\end{intuitionbox}

\begin{proposition}[Exact support-slack test]
\label{prop:support-slack-test}
For this proposition only, do not assume the standing inclusion
\(W(B)\subset\Omega\).  Let \(\Omega\) be a bounded convex domain with
\(C^1\) boundary, and assume
\(\partial\Omega\cap\spec(B)=\varnothing\).  At
\(\zeta\in\partial\Omega\), set \(u=\nu(\zeta)\) and
\[
  h_\Omega(u)
  =
  \max_{z\in\overline\Omega}\RePart(\overline uz)
  =
  \RePart(\overline u\zeta),
  \qquad
  D_{\Omega,B}(u)
  =
  h_\Omega(u)I-\RePart(\overline uB).
\]
Then the density in \eqref{eq:double-layer-measure} has the factorization
\begin{equation}
  d\mu_B(\zeta)
  =
  \frac1\pi
  \mathcal R_\zeta^*D_{\Omega,B}(u)\mathcal R_\zeta\,ds.
  \label{eq:support-slack-factorization}
\end{equation}
Moreover,
\[
  \frac{d\mu_B}{ds}(\zeta)\succeq0
  \text{ for every }\zeta\in\partial\Omega
  \quad\Longleftrightarrow\quad
  W(B)\subseteq\overline\Omega.
\]
\end{proposition}

\begin{proof}
The identity
\[
  \nu Z_\zeta^*+\overline\nu Z_\zeta
  =
  2D_{\Omega,B}(\nu)
\]
and congruence by the invertible resolvent give
\eqref{eq:support-slack-factorization}.  Thus positivity of the density
is equivalent to
\[
  \RePart(\overline uB)\preceq h_\Omega(u)I
\]
for every boundary normal \(u\).  By
\eqref{eq:W-support-function}, this says that every supporting
half-plane of \(\Omega\) contains \(W(B)\).  The intersection of those
half-planes is \(\overline\Omega\), proving the equivalence.
\end{proof}

\begin{remark}
If \(W(B)\) is compactly contained in \(\Omega\), then every
\(D_{\Omega,B}(u)\) is positive definite, uniformly in \(u\).
Moreover, since \(\mathcal R_\zeta\) is invertible,
\[
  \rank\left(\frac{d\mu_B}{ds}(\zeta)\right)
  =
  \rank D_{\Omega,B}(u).
\]
If \(h_\Omega(u)=h_B(u)\), where
\[
  h_B(u)
  =
  \lambda_{\max}\left(\RePart(\overline uB)\right),
\]
then
\[
  \ker D_{\Omega,B}(u)
  =
  \ker\left(
    h_B(u)I-\RePart(\overline uB)
  \right).
\]
Thus the nullity records the dimension of the supporting-state space.
The tangential compression from
\cref{prop:exposed-face-pencil} then determines the endpoints of the
exposed face.
\end{remark}

\section{The mass is \texorpdfstring{\(2I\)}{2I}}

Using \(d\zeta=i\nu\,ds\) and the matrix Cauchy formula,
\[
  \frac{1}{2\pi}\int_{\partial\Omega}
  \nu(\zeta)\mathcal R_\zeta\,ds
  =
  \frac{1}{2\pi i}\int_{\partial\Omega}
  (\zeta I-B)^{-1}\,d\zeta
  =I.
\]
Taking adjoints gives the same mass for the second layer.  Hence
\begin{equation}
  \int_{\partial\Omega}d\mu_B=2I.
  \label{eq:double-layer-mass}
\end{equation}
Define
\begin{equation}
  \Phi(\varphi)=\frac12
  \int_{\partial\Omega}\varphi(\zeta)\,d\mu_B(\zeta),
  \qquad \varphi\in C(\partial\Omega).
  \label{eq:Phi}
\end{equation}
By \eqref{eq:double-layer-mass}, \(\Phi(1)=I\); positivity of the
measure makes \(\Phi\) positive.  Thus \(\Phi\) is unital,
\(*\)-preserving, and contractive.

\begin{proposition}[Support--resolvent resolution of the identity]
\label{prop:support-resolvent-resolution}
Under the setup of this chapter,
\begin{equation}
  \int_{\partial\Omega}
  \mathcal R_\zeta^*
  D_{\Omega,B}(\nu(\zeta))
  \mathcal R_\zeta\,ds
  =2\pi I.
  \label{eq:support-resolvent-resolution}
\end{equation}
Equivalently, for every \(x\in\C^n\),
\begin{equation}
  \int_{\partial\Omega}
  \norm{
    D_{\Omega,B}(\nu(\zeta))^{1/2}
    \mathcal R_\zeta x
  }_2^2\,ds
  =2\pi\norm{x}_2^2.
  \label{eq:support-resolvent-vector}
\end{equation}
\end{proposition}

\begin{proof}
Insert the factorization
\eqref{eq:support-slack-factorization} into the mass identity
\eqref{eq:double-layer-mass}.
\end{proof}

\begin{intuitionbox}[A support-weighted conservation law]
Resolvent amplification is penalized by the support slack.  The total
slack-weighted boundary energy is exactly \(2\pi I\), independently of
the detailed shape of \(\Omega\).
\end{intuitionbox}

\section{The companion transform}

For \(h\) holomorphic on a neighborhood of
\(\overline\Omega\), define
\begin{equation}
  g_h(z)=\frac{1}{2\pi i}
  \int_{\partial\Omega}
  \frac{\overline{h(\zeta)}}{\zeta-z}\,d\zeta,
  \qquad z\in\Omega.
  \label{eq:companion}
\end{equation}
The boundary data \(\overline{h(\zeta)}\) are generally not
holomorphic in \(\zeta\), but the integral defines a holomorphic
function of the interior variable \(z\).

\begin{lemma}[Double-layer identity]
\label{lem:double-layer}
For \(h\) as above,
\begin{equation}
  2\Phi(h)=h(B)+g_h(B)^*.
  \label{eq:double-layer-identity}
\end{equation}
\end{lemma}

\begin{proof}
The first layer of \(2\Phi(h)\) is
\[
  \frac{1}{2\pi}
  \int_{\partial\Omega}
  h(\zeta)\nu(\zeta)(\zeta I-B)^{-1}\,ds
  =h(B)
\]
by \(d\zeta=i\nu ds\) and the holomorphic functional calculus.

For the second layer, do not assume that \(g_h\) extends across the
boundary.  Choose a contour \(\Gamma\subset\Omega\) enclosing
\(\spec(B)\).  Cauchy's formula and finite-dimensional Fubini give
\begin{align}
  g_h(B)
  &=\frac{1}{2\pi i}\int_\Gamma
    g_h(z)(zI-B)^{-1}\,dz \notag\\
  &=\frac{1}{2\pi i}\int_{\partial\Omega}
    \overline{h(\zeta)}(\zeta I-B)^{-1}\,d\zeta.
  \label{eq:companion-at-B}
\end{align}
Taking adjoints and using
\(d\overline\zeta=-i\overline\nu\,ds\) identifies this with the
second layer of \(2\Phi(h)\).
\end{proof}

\section{The boundary POVM and a symmetrized normal dilation}

The normalized measure
\[
  E_B(\mathcal E)=\frac12\mu_B(\mathcal E),
  \qquad \mathcal E\subseteq\partial\Omega\ \text{Borel},
\]
is a positive operator-valued measure with \(E_B(\partial\Omega)=I\);
in short, it is a POVM.  The map \(\Phi\) is integration against this
POVM.  Because its source \(C(\partial\Omega)\) is commutative, a
positive map from it is automatically completely positive.

\begin{blackbox}[Naimark dilation for the boundary POVM]
There are a Hilbert space \(\mathcal K\), an isometry
\(V:\C^n\to\mathcal K\), and a normal operator \(\mathcal N\) on
\(\mathcal K\), with
\(\spec(\mathcal N)\subseteq\partial\Omega\), such that
\[
  \Phi(\varphi)=V^*\varphi(\mathcal N)V
  \qquad
  (\varphi\in C(\partial\Omega)).
\]
This is the standard Naimark--Stinespring representation of a unital
positive map on a commutative \(C^*\)-algebra.  The dilation is not
unique, and it is not needed for the norm proof.
\end{blackbox}

In this representation, the double-layer identity becomes
\begin{equation}
  h(B)+g_h(B)^*
  =
  2V^*h(\mathcal N)V.
  \label{eq:symmetrized-normal-dilation}
\end{equation}
Thus numerical-range containment naturally dilates the
\emph{symmetrized} analytic calculus.  The obstruction is not the
absence of a positive boundary model; it is the need to remove the
coanalytic companion without treating it as an unrelated summand.

\begin{pitfall}
The identity \eqref{eq:double-layer-identity} is symmetrized.  A bound
on \(\Phi(h)\) controls the sum \(h(B)+g_h(B)^*\), not \(h(B)\) alone.
Estimating the companion as an unrelated summand throws away the
algebraic coupling that the new proof exploits.
\end{pitfall}

\section{What is classical and what will be new}

The positive boundary measure and companion transform are classical
ingredients of the Delyon--Delyon and Crouzeix--Palencia framework
\cite{DelyonDelyon1999,CrouzeixPalencia2017}.  The new move is to
apply \(\Phi\) not just to \(f\), but to the entire Cayley family
\[
  \frac{1+wf}{1-wf},
  \qquad w\in\D.
\]
This records all companion terms \(g_{f^m}(B)^*\) in one analytic
resolvent completion.

\section{Exercises}

\begin{exercise}[Warm-up: a scalar boundary point]
For \(\Omega=\D\), \(B=0\), and \(\zeta=e^{it}\), compute
\(\nu(\zeta)\), \(Z(\zeta)\), and the density
\(\overline\nu Z+\nu Z^*\) before reading off the mass of \(\Phi\).
\end{exercise}

\begin{exercise}[Core: support to density]
\label{ex:density-positive}
Write out the congruence that transforms \eqref{eq:support-matrix} into
the density in \eqref{eq:double-layer-measure}; verify the order of
all noncommuting factors.
\end{exercise}

\begin{exercise}[Core: two masses]
Derive \eqref{eq:double-layer-mass}, showing separately the first and
second layers.
\end{exercise}

\begin{exercise}[Core: companion identity]
\label{ex:companion-identity}
Fill in the Fubini calculation leading to
\eqref{eq:companion-at-B}.  Why is an interior contour \(\Gamma\)
important?
\end{exercise}

\begin{exercise}[Core: scalar disk]
Take \(n=1\), \(B=0\), and \(\Omega=\D\).  Compute the density
\eqref{eq:double-layer-measure} and the map \(\Phi\) explicitly.
\end{exercise}

\begin{exercise}[Stretch: no multiplicativity]
Find a simple positive unital map \(\Phi:C(X)\to\mathbb M_2(\C)\) and
a function \(f\) for which \(\Phi(f^2)\neq\Phi(f)^2\).
\end{exercise}

\part{The positive-real completion mechanism}

\chapter{The Cayley family creates a resolvent completion}
\label{ch:cayley-completion}

\begin{learninggoals}
\begin{itemize}
  \item Build a matrix-valued Carath\'eodory function from the
  double-layer map.
  \item Expand it coefficient by coefficient.
  \item Locate the resolvent defect in \(\alg(T^*)\).
  \item State the abstract positive-real completion theorem.
\end{itemize}
\end{learninggoals}

\begin{workshopplan}
\textbf{Opening:} revisit the scalar Cayley series.
\textbf{Core proof:} construct \(H\), subtract the resolvent, and
track the defect algebra.
\textbf{Error clinic:} why \(\Phi(f^m)\neq\Phi(f)^m\) is harmless.
\textbf{Studio:} scalar and \(2\times2\) completion examples.
\end{workshopplan}

\section{Auxiliary hypotheses}

For this week, assume:
\begin{itemize}
  \item \(B\in\Mn\);
  \item \(\Omega\) is a bounded convex \(C^1\) domain with
  \(W(B)\subset\Omega\);
  \item \(f\) is holomorphic on a neighborhood of
  \(\overline\Omega\) and
  \(\max_{\overline\Omega}|f|\leq1\);
  \item \(T=f(B)\) satisfies
  \begin{equation}
    \alg(T)=\alg(B).
    \label{eq:alg-generation}
  \end{equation}
\end{itemize}
The last condition is not automatic: a function can send different
eigenvalues of \(B\) to the same value.  Week 13 will enforce it by an
arbitrarily small perturbation.

By spectral mapping,
\[
  \spec(T)=f(\spec(B))\subset\olD.
\]

\section{Construction of the positive-real function}

For \(w\in\D\), define on the boundary
\begin{equation}
  c_w(\zeta)=\frac{1+wf(\zeta)}{1-wf(\zeta)}
  \label{eq:cayley-boundary}
\end{equation}
and put
\begin{equation}
  H(w)=\Phi(c_w),
  \label{eq:H-Phi}
\end{equation}
where \(\Phi\) is the double-layer map from Week 8.

The map \(w\mapsto c_w\) is analytic with values in
\(C(\partial\Omega)\), locally uniformly on \(\D\).  Since \(\Phi\)
is bounded and linear, \(H\) is analytic.  Also
\[
  H(0)=\Phi(1)=I.
\]
By \eqref{eq:cayley-real},
\[
  \RePart c_w(\zeta)
  =\frac{1-|wf(\zeta)|^2}{|1-wf(\zeta)|^2}\geq0.
\]
Because \(\Phi\) is positive and \(*\)-preserving,
\begin{equation}
  \RePart H(w)
  =\Phi(\RePart c_w)\succeq0.
  \label{eq:H-positive}
\end{equation}
Thus \(H\) is a normalized matrix-valued
\emph{Carath\'eodory function}: analytic on \(\D\), equal to \(I\) at
the origin, with positive semidefinite real part.

\section{Coefficient ledger}

The boundary expansion
\[
  c_w=1+2\sum_{m=1}^{\infty}w^mf^m
\]
converges in the supremum norm on \(\partial\Omega\), locally uniformly
for \(w\in\D\).  Apply \(\Phi\) term by term:
\[
  H(w)=I+2\sum_{m=1}^{\infty}w^m\Phi(f^m).
\]
The double-layer identity, applied separately to each \(f^m\), says
\[
  2\Phi(f^m)=f(B)^m+g_{f^m}(B)^*
  =T^m+g_{f^m}(B)^*.
\]
Therefore
\begin{equation}
  H(w)=I+\sum_{m=1}^{\infty}w^m
  \bigl(T^m+g_{f^m}(B)^*\bigr).
  \label{eq:H-series}
\end{equation}

\begin{checkpoint}[Why the coefficients are exactly right]
The Cayley series contributes \(2\); the double-layer identity divides
that coefficient between \(T^m\) and its companion.  The known
coefficients are \(T^m\), so their sum is the resolvent:
\[
  I+\sum_{m\geq1}w^mT^m=(I-wT)^{-1}.
\]
\end{checkpoint}

\section{The defect lies in the adjoint algebra}

Because \(\spec(T)\subset\olD\) and \(w\in\D\), the resolvent series
converges in norm.  Subtracting it from \eqref{eq:H-series} gives
\begin{equation}
  H(w)-(I-wT)^{-1}
  =\sum_{m=1}^{\infty}w^m g_{f^m}(B)^*.
  \label{eq:completion-series}
\end{equation}
Each \(g_{f^m}\) is holomorphic near \(\spec(B)\).  By finite-dimensional
holomorphic functional calculus,
\[
  g_{f^m}(B)\in\alg(B).
\]
Taking adjoints and using \eqref{eq:alg-generation},
\[
  g_{f^m}(B)^*
  \in\alg(B^*)=\alg(T^*).
\]
Every partial sum in \eqref{eq:completion-series} lies in
\(\alg(T^*)\), and this finite-dimensional algebra is closed.
Consequently,
\begin{equation}
  H(w)-(I-wT)^{-1}\in\alg(T^*)
  \qquad(w\in\D).
  \label{eq:completion-membership}
\end{equation}

\begin{pitfall}[Three different functions]
\(c_w\) is a scalar boundary function of \(\zeta\);
\(H(w)\) is a matrix function of \(w\); and
\((I-wT)^{-1}\) is the known matrix resolvent.  Keeping these roles
separate prevents the false identity
\(H(w)=(I+wT)(I-wT)^{-1}\).
\end{pitfall}

\section{The abstract theorem now needed}

\begin{theorem}[Positive-real completion]
\label{thm:completion}
Let \(T\in\Mn\) have \(n\) distinct eigenvalues, all in \(\olD\).
Suppose \(H:\D\to\Mn\) is analytic and satisfies
\begin{align}
  H(0)&=I,
  &
  \RePart H(w)&\succeq0
  \quad(w\in\D),
  \label{eq:completion-positive}\\
  H(w)-(I-wT)^{-1}&\in\alg(T^*)
  \quad(w\in\D).
  \label{eq:completion-algebra}
\end{align}
Then
\[
  \norm{T}\leq2.
\]
\end{theorem}

This theorem is independent of numerical ranges and double-layer
potentials.  It is the new abstract tool.  Weeks 10--12 prove it.

\begin{corollary}[Auxiliary sharp bound]
\label{cor:auxiliary}
Under the hypotheses of this chapter, if \(T=f(B)\) has distinct
eigenvalues, then \(\norm{f(B)}\leq2\).
\end{corollary}

\begin{proof}
The construction above verifies all hypotheses of
\cref{thm:completion}.
\end{proof}

\section{Why the full family matters}

The single identity
\[
  2\Phi(f)=T+g_f(B)^*
\]
contains one companion matrix.  The full family
\eqref{eq:completion-series} says much more:
\begin{itemize}
  \item the known part is an entire resolvent;
  \item every unknown coefficient is adjoint-aligned;
  \item the unknowns form one analytic defect in a fixed finite
  algebra.
\end{itemize}
The completion theorem uses precisely this extra structure.

\begin{remark}[Adjoint-algebraic unitary moment completion]
\label{rem:unitary-moment-completion}
Put \(C_m=g_{f^m}(B)\).  Comparing \eqref{eq:H-series} with the matrix
Herglotz representation gives a positive \(\Mn\)-valued measure \(M\)
on \(\Torus\).  Because \(H(0)=I\), the imaginary constant in the
Herglotz representation vanishes and \(M(\Torus)=I\).  Moreover,
\[
  T^m+C_m^*
  =
  2\int_{\Torus}\overline\zeta^{\,m}\,dM(\zeta)
  \qquad(m\geq1).
\]
Let \(F\) be a Naimark dilation of \(M\), with isometry \(J\), and
define the unitary
\[
  U=\int_{\Torus}\overline\zeta\,dF(\zeta).
\]
Then
\begin{equation}
  T^m+C_m^*=2J^*U^mJ
  \qquad(m\geq1).
  \label{eq:unitary-moment-completion}
\end{equation}
Under the algebra-generation hypothesis, every \(C_m^*\) belongs to
\(\alg(T^*)\).  Thus the displayed positive-index sequence extends,
by taking zeroth term \(2I\) and negative terms as adjoints, to a
positive-definite two-sided unitary moment sequence.  The completion
terms remain inside the opposite functional calculus.

For \(\Omega=\D\) and \(f(z)=z\), the boundary calculation gives
\(g_{z^m}=0\) for every \(m\geq1\).  In that case
\eqref{eq:unitary-moment-completion} reduces to the unitary
\(2\)-dilation identity \(B^m=2J^*U^mJ\).  On a general convex domain,
the companion \(C_m^*\) measures the failure of this disk
Fourier splitting, while the proof shows that the failure remains
algebraically constrained rather than arbitrary.
\end{remark}

\section{Exercises}

\begin{exercise}[Warm-up: Cayley coefficients]
Expand \((1+wa)/(1-wa)\) through order \(w^3\) and identify the
coefficient of \(w^m\) for \(m\geq1\).
\end{exercise}

\begin{exercise}[Core: positivity]
Starting from \eqref{eq:cayley-boundary}, prove
\eqref{eq:H-positive} without assuming that \(\Phi\) is multiplicative.
\end{exercise}

\begin{exercise}[Core: coefficient audit]
\label{ex:coefficient-audit}
Derive \eqref{eq:H-series} from the Cayley and double-layer formulas,
showing where both factors \(2\) go.
\end{exercise}

\begin{exercise}[Core: algebra audit]
Explain why Hermite rather than merely Lagrange interpolation may be
needed to prove \(g_{f^m}(B)\in\alg(B)\) at this stage.
\end{exercise}

\begin{exercise}[Core: scalar completion]
When \(n=1\), write \cref{thm:completion} in scalar language.  What does
the defect condition say, and why is the conclusion unsurprising?
\end{exercise}

\begin{exercise}[Stretch: a wrong orientation]
\label{ex:wrong-orientation}
Let
\[
  T_M=\begin{pmatrix}0&M\\0&1/2\end{pmatrix},
  \qquad H(w)\equiv I,
  \qquad M>0.
\]
Show that
\[
  H(w)-(I-wT_M)^{-1}\in\alg(T_M)
\]
for all \(w\in\D\), while \(\norm{T_M}\to\infty\) as \(M\to\infty\).
Thus \(\alg(T^*)\) cannot be replaced by \(\alg(T)\).
\end{exercise}

\chapter{Diagonalizing the structured nuisance term}
\label{ch:diagonal-nuisance}

\begin{learninggoals}
\begin{itemize}
  \item Convert the adjoint-algebra defect into a diagonal analytic
  correction.
  \item Preserve positivity using congruence.
  \item Understand the Gram matrix \(G=S^*S\).
  \item Set up the special kernel samples.
\end{itemize}
\end{learninggoals}

\begin{workshopplan}
\textbf{Opening:} diagonalize one nonnormal \(2\times2\) matrix and
compute its Gram matrix. \textbf{Core:} pull the defect back by
congruence and track the order \(\Theta G\).
\textbf{Guided proof:} derive the half-eigenvalue and origin samples.
\textbf{Studio:} reproduce the cancellation row by row, then complete
the change-of-basis audit.
\end{workshopplan}

\section{Diagonalization and its Gram matrix}

Assume the hypotheses of \cref{thm:completion}.  Since \(T\) has
distinct eigenvalues,
\begin{equation}
  T=S\Lambda S^{-1},
  \qquad
  \Lambda=\diag(\lambda_1,\ldots,\lambda_n),
  \qquad
  G=S^*S.
  \label{eq:T-diagonalization}
\end{equation}
The eigenvalues lie in \(\olD\), and \(G\succ0\).
The matrix \(G\) is the Gram matrix of the possibly nonorthogonal
eigenbasis.  It stores the nonnormal geometry that \(\Lambda\) alone
cannot see.

By \cref{lem:alg-simple},
\begin{equation}
  \alg(T^*)=
  \set{(S^{-1})^*DS^*:D\text{ diagonal}}.
  \label{eq:adjoint-diagonal-algebra}
\end{equation}

\section{The defect becomes diagonal}

Let
\[
  \Delta(w)=H(w)-(I-wT)^{-1}.
\]
For each \(w\), membership in \eqref{eq:adjoint-diagonal-algebra}
gives a unique diagonal matrix \(\Theta(w)\) such that
\[
  \Delta(w)=(S^{-1})^*\Theta(w)S^*.
\]
Explicitly,
\[
  \Theta(w)=S^*\Delta(w)(S^*)^{-1}.
\]
Therefore \(\Theta\) is analytic.  Since \(H(0)=I\),
\(\Delta(0)=0\), so
\begin{equation}
  \Theta(0)=0.
  \label{eq:Theta-zero}
\end{equation}
Write
\[
  \Theta(w)=\diag(\theta_1(w),\ldots,\theta_n(w)).
\]

Now pull \(H\) back by congruence:
\[
  \widetilde H(w)=S^*H(w)S.
\]
The resolvent part transforms as
\[
  S^*(I-wT)^{-1}S
  =G(I-w\Lambda)^{-1},
\]
while the defect becomes \(\Theta(w)G\).  Hence
\begin{equation}
  \widetilde H(w)
  =G(I-w\Lambda)^{-1}+\Theta(w)G.
  \label{eq:Htilde}
\end{equation}

\begin{pitfall}[Factor order matters]
The nuisance term is \(\Theta(w)G\), not \(G\Theta(w)\).  These
matrices need not commute.  The rowwise form \(\Theta G\) is exactly
what the sampling vectors can annihilate.
\end{pitfall}

\section{Positivity survives the change of coordinates}

Similarity does not preserve positive real part, but congruence does:
\[
  \RePart\widetilde H(w)
  =S^*(\RePart H(w))S\succeq0.
\]
Thus the Herglotz kernel
\[
  \mathcal L_{\widetilde H}(w,z)
  =\frac{\widetilde H(w)+\widetilde H(z)^*}
         {1-w\overline z}
\]
is positive by \cref{lem:herglotz-kernel}.

\section{The two Cauchy-type matrices}

Define Hermitian matrices \(P,Q,Y\) by
\begin{equation}
  P_{ij}=\frac{G_{ij}}
  {1-\overline{\lambda_i}\lambda_j/4},
  \qquad
  Q_{ij}=\frac{G_{ij}}
  {1-\overline{\lambda_i}\lambda_j/2},
  \qquad
  Y=Q-P.
  \label{eq:PQY}
\end{equation}
The denominators are nonzero because
\(|\overline{\lambda_i}\lambda_j|\leq1\).  Hermitian symmetry follows
by conjugating the \(ij\)-entry.

At this point these definitions may look engineered.  They will arise
automatically from the two denominators in the sampled Herglotz
kernel:
\[
  1-w_i\overline{w_j}
  \quad\text{and}\quad
  1-w_i\lambda_j.
\]

\section{Choose the sample points and directions}

Fix an arbitrary vector \(u=(u_1,\ldots,u_n)^T\).  For
\(1\leq i\leq n\), take
\begin{equation}
  w_i=\frac{\overline{\lambda_i}}2,
  \qquad
  \xi_i=u_i e_i.
  \label{eq:eigen-samples}
\end{equation}
Even when \(|\lambda_i|=1\), the sample lies strictly inside \(\D\).
Add one more sample at the origin:
\begin{equation}
  w_0=0,
  \qquad
  \xi_0=v=-G^{-1}Pu.
  \label{eq:origin-sample}
\end{equation}
The origin vector is allowed to depend on \(u\).  Kernel positivity
holds for every finite collection of vectors, so this is legitimate.

\section{Exact cancellation of the unknown function}

Let \(E(w)=\Theta(w)G\) be the nuisance part of
\eqref{eq:Htilde}.  Because \(\Theta(0)=0\), all terms in the sampled
kernel quadratic form involving \(E(0)\) vanish.  The remaining
contribution of \(E\) and \(E^*\) is twice the real part of
\begin{equation}
  \sum_{i=1}^n\overline{u_i}\theta_i(w_i)
  \left(
    (Gv)_i+
    \sum_{j=1}^n
    \frac{G_{ij}u_j}{1-w_i\overline{w_j}}
  \right).
  \label{eq:theta-contribution}
\end{equation}
Here is the index calculation:
\[
  e_i^*\Theta(w_i)G=\theta_i(w_i)e_i^*G,
\]
so the diagonal correction selects row \(i\) of \(G\).  Moreover,
\[
  1-w_i\overline{w_j}
  =1-\frac{\overline{\lambda_i}\lambda_j}{4}.
\]
Therefore the sum in the large parentheses is \((Pu)_i\), and the
whole parenthesis in \eqref{eq:theta-contribution} equals
\[
  (Gv+Pu)_i.
\]
Our definition \(v=-G^{-1}Pu\) makes this zero for every \(i\).
The entire unknown analytic correction disappears exactly.

\begin{proofmap}[Why one origin vector is enough]
\[
\begin{array}{c}
\Delta(w)\in\alg(T^*)\\
\Downarrow\\
\Theta(w)\text{ is diagonal}\\
\Downarrow\\
n\text{ unknown row directions}\\
\Downarrow\\
Gv+Pu=0\text{ is a system of exactly \(n\) equations}\\
\Downarrow\\
v=-G^{-1}Pu\text{ cancels all rows simultaneously.}
\end{array}
\]
An arbitrary matrix-valued correction would have \(n^2\) independent
directions and would not be removable this way.
\end{proofmap}

\begin{pitfall}[Cancellation is not vanishing]
We have not proved \(\Theta\equiv0\).  We have only shown that
\(\Theta\)'s contribution to one specially designed family of kernel
quadratic forms is zero.  That is all the proof needs.
\end{pitfall}

\section{Exercises}

\begin{exercise}[Warm-up: one row of a Gram matrix]
For
\(
S=\begin{pmatrix}1&2\\0&1\end{pmatrix}
\),
compute \(G=S^*S\), its inverse, and its two rows.  Explain why
\(\Theta G\) scales rows whereas \(G\Theta\) scales columns.
\end{exercise}

\begin{exercise}[Core: transform the defect]
Derive \eqref{eq:Htilde} from
\eqref{eq:T-diagonalization} and
\eqref{eq:adjoint-diagonal-algebra}, showing every multiplication by
\(S\) or \(S^*\).
\end{exercise}

\begin{exercise}[Core: analyticity and normalization]
Prove that \(\Theta\) is analytic and that
\eqref{eq:Theta-zero} follows from \(H(0)=I\).
\end{exercise}

\begin{exercise}[Core: correction calculation]
\label{ex:theta-cancel}
Starting from the definition of the Herglotz kernel, derive
\eqref{eq:theta-contribution} without skipping the adjoint half.
\end{exercise}

\begin{exercise}[Core: dimension count]
Explain carefully why diagonal structure reduces the correction from
\(n^2\) scalar directions to \(n\) row directions.
\end{exercise}

\begin{exercise}[Stretch: \(2\times2\) symbolic check]
Take arbitrary distinct
\(\lambda_1,\lambda_2\in\olD\), a positive matrix
\(G=\begin{pmatrix}a&b\\\overline b&d\end{pmatrix}\), and diagonal
\(\Theta\).  Write \(P\), \(v\), and the two cancellation equations
explicitly.
\end{exercise}

\chapter{The surviving kernel inequality and two Gramians}
\label{ch:gramian-inequality}

\begin{learninggoals}
\begin{itemize}
  \item Compute every block of the sampled known kernel.
  \item Derive the anticommutator inequality for \(Y=Q-P\).
  \item Balance the eigenbasis.
  \item Identify \(P\) and \(Q\) as weighted Gramians.
\end{itemize}
\end{learninggoals}

\begin{workshopplan}
\textbf{Opening:} verify the scalar partial fraction that creates
\(P\) and \(Q\). \textbf{Core:} assemble the three known kernel
blocks. \textbf{Guided proof:} substitute the compensating vector and
balance the basis. \textbf{Studio:} trace the worked \(2\times2\)
audit and then repeat it with new eigenvalues.
\end{workshopplan}

\section{The sample--sample block}

The known part of \(\widetilde H\) is
\[
  F(w)=G(I-w\Lambda)^{-1}.
\]
At sample points \(w_i=\overline{\lambda_i}/2\), the \(ij\)-entry of
its Herglotz kernel is
\begin{align}
  \frac{F(w_i)_{ij}+F(w_j)^*_{ij}}
       {1-w_i\overline{w_j}}
  &=
  \frac{2G_{ij}}
  {(1-\overline{\lambda_i}\lambda_j/2)
   (1-\overline{\lambda_i}\lambda_j/4)}
  \notag\\
  &=(4Q-2P)_{ij}.
  \label{eq:main-sample-block}
\end{align}
The last equality is the scalar partial-fraction identity
\[
  \frac{2}{(1-x/2)(1-x/4)}
  =\frac{4}{1-x/2}-\frac{2}{1-x/4}.
\]

\section{The origin blocks}

For a fixed eigenvalue sample, the literal sample--origin block is
\[
  \mathcal L_F(w_i,0)=F(w_i)+G.
\]
After applying the prescribed vector \(\xi_i=u_i e_i\), its selected
row is
\[
  e_i^*\bigl(F(w_i)+G\bigr)=e_i^*(Q+G).
\]
Thus, after the \(n\) selected rows are assembled in the quadratic
form, the \emph{effective} cross block is
\begin{equation}
  G+Q.
  \label{eq:cross-block}
\end{equation}
At the origin the literal and effective block is
\begin{equation}
  2G.
  \label{eq:origin-block}
\end{equation}

After the exact cancellation in Week 10, positivity of the full
Herglotz kernel therefore yields, for the particular
\(v=-G^{-1}Pu\),
\begin{equation}
  0\leq
  \begin{pmatrix}u\\v\end{pmatrix}^{\!*}
  \begin{pmatrix}
    4Q-2P&G+Q\\
    G+Q&2G
  \end{pmatrix}
  \begin{pmatrix}u\\v\end{pmatrix}.
  \label{eq:block-quadratic}
\end{equation}

\begin{pitfall}
Equation \eqref{eq:block-quadratic} is asserted on the graph
\(v=-G^{-1}Pu\).  The displayed known-part block matrix is not claimed
to be positive semidefinite for arbitrary independent \(u,v\); the
unknown part was canceled only for the special graph.
\end{pitfall}

\section{Substitute the compensator}

Insert \(v=-G^{-1}Pu\) into \eqref{eq:block-quadratic}.  Since \(P\),
\(Q\), and \(G\) are Hermitian, the result is
\begin{align}
0\leq u^*\bigl[
 &(4Q-2P)
 -(G+Q)G^{-1}P
 -PG^{-1}(G+Q)
 +2PG^{-1}P
\bigr]u.
\label{eq:before-Y}
\end{align}
Use \(Q=P+Y\).  The pure \(P\)-terms cancel:
\begin{align*}
&(4Q-2P)
 -(G+Q)G^{-1}P
 -PG^{-1}(G+Q)
 +2PG^{-1}P\\
&\qquad
=4Y-YG^{-1}P-PG^{-1}Y.
\end{align*}
Since \(u\) was arbitrary,
\begin{equation}
  4Y-YG^{-1}P-PG^{-1}Y\succeq0.
  \label{eq:Y-inequality}
\end{equation}
This is the entire payload of the kernel sampling.

\section{A fully worked \texorpdfstring{\(2\times2\)}{2 by 2} audit}

The following small case makes the cancellation visible without
hiding it in indices.  Take
\[
  T=\begin{pmatrix}0&1\\0&1/2\end{pmatrix}
  =S\Lambda S^{-1},
  \qquad
  S=\begin{pmatrix}1&2\\0&1\end{pmatrix},
  \qquad
  \Lambda=\diag(0,1/2).
\]
Then
\[
  G=S^*S=\begin{pmatrix}1&2\\2&5\end{pmatrix},
  \qquad
  G^{-1}=\begin{pmatrix}5&-2\\-2&1\end{pmatrix}.
\]
The definitions \eqref{eq:PQY} give
\[
  P=\begin{pmatrix}1&2\\2&16/3\end{pmatrix},
  \qquad
  Q=\begin{pmatrix}1&2\\2&40/7\end{pmatrix},
  \qquad
  Y=\begin{pmatrix}0&0\\0&8/21\end{pmatrix}.
\]
For \(u=(u_1,u_2)^T\),
\[
  G^{-1}P=
  \begin{pmatrix}1&-2/3\\0&4/3\end{pmatrix},
  \qquad
  v=-G^{-1}Pu
  =\begin{pmatrix}-u_1+\tfrac23u_2\\-\tfrac43u_2\end{pmatrix}.
\]
Consequently \(Gv+Pu=0\).  Whatever the two analytic diagonal entries
of \(\Theta\) may be, their sampled contribution
\[
  2\RePart\sum_{i=1}^2
  \overline{u_i}\theta_i(w_i)(Gv+Pu)_i
\]
is exactly zero.  The known remainder is positive because
\[
  4Y-YG^{-1}P-PG^{-1}Y
  =\begin{pmatrix}0&0\\0&32/63\end{pmatrix}\succeq0.
\]
This calculation is an audit of the completion mechanism: it assumes
the positive-real completion hypotheses, but it does not need to know
the values of the nuisance functions \(\theta_i\).

\section{Balance the eigenbasis}

Define
\begin{equation}
\begin{aligned}
  \widetilde T&=G^{1/2}\Lambda G^{-1/2},&
  \widehat P&=G^{-1/2}PG^{-1/2},\\
  \widehat Q&=G^{-1/2}QG^{-1/2},&
  \widehat Y&=\widehat Q-\widehat P.
\end{aligned}
\label{eq:balanced}
\end{equation}
This is not an arbitrary normalization.  The polar decomposition of
\(S\) is
\[
  S=UG^{1/2}
\]
with \(U\) unitary, and therefore
\begin{equation}
  T=U\widetilde T U^*.
  \label{eq:T-unitary-tilde}
\end{equation}
Thus \(T\) and \(\widetilde T\) have the same operator norm.

Congruence of \eqref{eq:Y-inequality} by \(G^{-1/2}\) gives
\begin{equation}
  4\widehat Y-\widehat Y\widehat P
  -\widehat P\widehat Y\succeq0.
  \label{eq:Yhat-inequality}
\end{equation}

\section{The two weighted Gramians}

Expand the scalar denominators in \eqref{eq:PQY} as geometric series.
Entrywise,
\[
  P=\sum_{k=0}^{\infty}
  4^{-k}\Lambda^{*k}G\Lambda^k,
  \qquad
  Q=\sum_{k=0}^{\infty}
  2^{-k}\Lambda^{*k}G\Lambda^k.
\]
After congruence,
\begin{equation}
  \widehat P=
  \sum_{k=0}^{\infty}
  4^{-k}\widetilde T^{*k}\widetilde T^k,
  \qquad
  \widehat Q=
  \sum_{k=0}^{\infty}
  2^{-k}\widetilde T^{*k}\widetilde T^k.
  \label{eq:two-gramians}
\end{equation}
These are weighted observability Gramians: on a vector \(x\),
\[
  x^*\widehat P x
  =\sum_{k\geq0}4^{-k}\norm{\widetilde T^kx}_2^2.
\]

Even if some \(|\lambda_j|=1\), the series converge in norm.  Indeed,
\[
  \widetilde T^k=G^{1/2}\Lambda^kG^{-1/2},
\]
so \(\{\widetilde T^k\}\) is uniformly bounded; the scalar weights are
summable.

Subtracting the two series gives
\begin{equation}
  \widehat Y=
  \sum_{k=1}^{\infty}
  (2^{-k}-4^{-k})
  \widetilde T^{*k}\widetilde T^k
  \succeq0.
  \label{eq:Yhat-positive}
\end{equation}
Also, the \(k=0\) term in \(\widehat P\) gives
\begin{equation}
  \widehat P\succeq I.
  \label{eq:Phat-lower}
\end{equation}

\begin{intuitionbox}[What the two scales accomplish]
The \(1/2\)-weighted Gramian \(\widehat Q\) sees more of the future
orbit than the \(1/4\)-weighted Gramian \(\widehat P\).
Their difference \(\widehat Y\) is a positive detector: if it sees
zero energy, then already \(\widetilde T x=0\).  The narrower Gramian
\(\widehat P\) also obeys the Stein recursion that will control
\(\norm{\widetilde T}\).
\end{intuitionbox}

\section{Exercises}

\begin{exercise}[Warm-up: the two denominators]
With \(x=1/2\), evaluate
\(
2/[(1-x/2)(1-x/4)]
\)
directly and from
\(
4/(1-x/2)-2/(1-x/4)
\).
\end{exercise}

\begin{exercise}[Core: top-left block]
\label{ex:top-left-block}
Derive \eqref{eq:main-sample-block} entry by entry.  Check both the
resolvent denominator and the Herglotz denominator.
\end{exercise}

\begin{exercise}[Core: graph substitution]
Starting with \eqref{eq:block-quadratic}, derive
\eqref{eq:before-Y} and then \eqref{eq:Y-inequality}.  Preserve the
order of all factors.
\end{exercise}

\begin{exercise}[Core: balance]
Prove \eqref{eq:T-unitary-tilde} from the polar decomposition of \(S\).
\end{exercise}

\begin{exercise}[Core: Gramian expansion]
\label{ex:gramian-expansion}
Prove both identities in \eqref{eq:two-gramians}, starting from the
entrywise formulas for \(P,Q\).
\end{exercise}

\begin{exercise}[Stretch: kernel of the detector]
Use \eqref{eq:Yhat-positive} to prove
\[
  \ker\widehat Y\subseteq\ker\widetilde T.
\]
Then show
\(
\ker\widehat Y\subseteq\ker(\widehat P-I).
\)
\end{exercise}

\chapter{Gramian rigidity, the Stein identity, and the constant
\texorpdfstring{\(2\)}{2}}
\label{ch:completion-finish}

\begin{learninggoals}
\begin{itemize}
  \item Convert an anticommutator inequality into a spectral bound.
  \item Use a positive detector even when it is singular.
  \item Derive and exploit a Stein identity.
  \item Explain conceptually why the sampling radius \(1/2\) is
  optimal within the one-parameter family studied below.
  \item Interpret the two factors \(\sqrt2\) as a controlled
  similarity and a transformed norm.
\end{itemize}
\end{learninggoals}

\begin{workshopplan}
\textbf{Opening:} test the anticommutator inequality on eigenvectors
of \(P\). \textbf{Core:} detector rigidity and its singular case.
\textbf{Guided proof:} shift the Gramian series to obtain the Stein
identity. \textbf{Oral studio:} reconstruct the completion theorem
from the seven-item checklist.
\end{workshopplan}

\section{A reusable rigidity lemma}

\begin{lemma}[Detector--anticommutator rigidity]
\label{lem:detector-rigidity}
Let \(P=P^*\) and \(X=X^*\) satisfy
\begin{equation}
  P\succeq I,\qquad
  X\succeq0,\qquad
  cX-XP-PX\succeq0,
  \qquad c\geq2.
  \label{eq:rigidity-hypotheses}
\end{equation}
Assume also the detectability condition
\begin{equation}
  \ker X\subseteq\ker(P-I).
  \label{eq:detectability}
\end{equation}
Then
\[
  P\preceq\frac c2 I.
\]
\end{lemma}

\begin{proof}
Choose a unit eigenvector \(e\) of the Hermitian matrix \(P\), with
\(Pe=\alpha e\).  Testing the anticommutator inequality on \(e\)
gives
\[
  0\leq(c-2\alpha)e^*Xe.
\]
If \(\alpha>c/2\), the coefficient is negative.  Since \(X\succeq0\),
we must have \(e^*Xe=0\).  By \cref{lem:psd-zero}, \(Xe=0\), and
detectability gives \(Pe=e\).  This contradicts
\(\alpha>c/2\geq1\).  Hence every eigenvalue of \(P\) is at most
\(c/2\).
\end{proof}

\begin{pitfall}
The anticommutator inequality alone is not enough when \(X\) is
singular.  The detector condition \eqref{eq:detectability} supplies
the missing rigidity on \(\ker X\).
\end{pitfall}

\section{Apply rigidity to the two Gramians}

Take
\[
  P=\widehat P,\qquad X=\widehat Y,\qquad c=4.
\]
The lower bound \(P\succeq I\) is \eqref{eq:Phat-lower};
positivity \(X\succeq0\) is \eqref{eq:Yhat-positive}; and the
anticommutator inequality is \eqref{eq:Yhat-inequality}.

Detectability follows from the series.  If
\(\widehat Y e=0\), then \(e^*\widehat Ye=0\), and
\eqref{eq:Yhat-positive} gives
\[
  0=\sum_{k=1}^{\infty}
  (2^{-k}-4^{-k})\norm{\widetilde T^ke}_2^2.
\]
Every coefficient is positive, so \(\widetilde Te=0\).  The series
for \(\widehat P\) then yields \(\widehat Pe=e\).
Thus \eqref{eq:detectability} holds, and
\cref{lem:detector-rigidity} gives
\begin{equation}
  I\preceq\widehat P\preceq2I.
  \label{eq:Phat-bounds}
\end{equation}

\section{The Stein identity}

Shift the first Gramian series:
\begin{align*}
  \widehat P-I
  &=\sum_{k=1}^{\infty}
    4^{-k}\widetilde T^{*k}\widetilde T^k\\
  &=\frac14\widetilde T^*
    \left(
      \sum_{j=0}^{\infty}
      4^{-j}\widetilde T^{*j}\widetilde T^j
    \right)\widetilde T.
\end{align*}
Hence
\begin{equation}
  \widehat P
  =I+\frac14\widetilde T^*\widehat P\widetilde T.
  \label{eq:stein}
\end{equation}
Using \eqref{eq:Phat-bounds},
\[
  \widetilde T^*\widetilde T
  \preceq
  \widetilde T^*\widehat P\widetilde T
  =4(\widehat P-I)
  \preceq4I.
\]
Therefore
\[
  \norm{\widetilde T}\leq2.
\]
Finally \(T=U\widetilde T U^*\) by
\eqref{eq:T-unitary-tilde}, so
\[
  \boxed{\norm{T}\leq2.}
\]
This completes the proof of \cref{thm:completion}.

\begin{checkpoint}[The sharp chain]
\[
\begin{aligned}
\text{positive kernel}
&\Longrightarrow
4\widehat Y-\widehat Y\widehat P-\widehat P\widehat Y\succeq0,\\
\widehat Y\text{ detects }\widetilde T
&\Longrightarrow I\preceq\widehat P\preceq2I,\\
\widehat P=I+\tfrac14\widetilde T^*\widehat P\widetilde T
&\Longrightarrow \widetilde T^*\widetilde T\preceq4I.
\end{aligned}
\]
The constant \(2\) is the square root of the final coefficient \(4\).
\end{checkpoint}

\begin{remark}[Equality cannot occur at the auxiliary stage]
\label{rem:auxiliary-strict}
If \(\norm{\widetilde T}=2\), choose a unit norm-attaining vector \(e\).
Equality throughout the Stein chain forces
\[
  \widehat Pe=2e,
  \qquad
  \widehat P\widetilde Te=\widetilde Te.
\]
The series for \(\widehat P-I\) then gives
\(\widetilde T^2e=0\), while
\(\norm{\widetilde Te}=2\).  This is a nontrivial Jordan chain at
zero.  But the completion theorem assumes simple spectrum, hence
diagonalizability, and a diagonalizable matrix satisfies
\(\ker\widetilde T^2=\ker\widetilde T\).  Therefore the auxiliary
estimate is actually strict:
\[
  \norm{T}<2.
\]
Accordingly, any equality obtained after the approximation steps must
emerge as diagonalizability is lost in the limiting argument.
\end{remark}

\section{Why one half is the distinguished radius}
\label{sec:general-radius}

The source proof samples at half the conjugate eigenvalues.  A
one-parameter version explains this choice.  Fix \(0<r<1\) and instead
sample
\[
  w_i=r\overline{\lambda_i}.
\]
Define
\begin{align}
  (P_r)_{ij}
  &=\frac{G_{ij}}{1-r^2\overline{\lambda_i}\lambda_j},
  &
  (Q_r)_{ij}
  &=\frac{G_{ij}}{1-r\overline{\lambda_i}\lambda_j},
  &
  Y_r&=Q_r-P_r.
  \label{eq:general-PQ}
\end{align}
The same origin choice \(v=-G^{-1}P_ru\) cancels the diagonal
correction.  The sample--sample block uses
\begin{equation}
  \frac{2}{(1-rx)(1-r^2x)}
  =
  \frac{2}{1-r}\frac{1}{1-rx}
  -\frac{2r}{1-r}\frac{1}{1-r^2x}.
  \label{eq:general-partial-fraction}
\end{equation}
Thus the three assembled known blocks are
\[
  \frac{2}{1-r}Q_r-\frac{2r}{1-r}P_r,
  \qquad
  G+Q_r,
  \qquad
  2G,
\]
in the sample--sample, cross, and origin positions, respectively.
Substitute \(v=-G^{-1}P_ru\) into their quadratic form and use
\(Q_r=P_r+Y_r\).  The terms containing only \(P_r\) cancel, leaving
\begin{equation}
  \frac{2}{1-r}Y_r
  -Y_rG^{-1}P_r
  -P_rG^{-1}Y_r\succeq0.
  \label{eq:general-unbalanced}
\end{equation}
Define the balanced matrices
\[
  \widehat P_r=G^{-1/2}P_rG^{-1/2},
  \qquad
  \widehat Y_r=G^{-1/2}Y_rG^{-1/2}.
\]
Congruence of \eqref{eq:general-unbalanced} gives
\begin{equation}
  \frac{2}{1-r}\widehat Y_r
  -\widehat Y_r\widehat P_r
  -\widehat P_r\widehat Y_r\succeq0,
  \label{eq:general-anticommutator}
\end{equation}
where
\begin{align}
  \widehat P_r
  &=\sum_{k=0}^{\infty}
    r^{2k}\widetilde T^{*k}\widetilde T^k,
  \label{eq:Pr-series}\\
  \widehat Y_r
  &=\sum_{k=1}^{\infty}
    (r^k-r^{2k})\widetilde T^{*k}\widetilde T^k
    \succeq0.
  \label{eq:Yr-series}
\end{align}
The same detector argument applies with
\(
c=2/(1-r)
\),
giving
\[
  I\preceq\widehat P_r\preceq\frac{1}{1-r}I.
\]
The Stein identity now reads
\[
  \widehat P_r
  =I+r^2\widetilde T^*\widehat P_r\widetilde T.
\]
Therefore
\[
  r^2\widetilde T^*\widetilde T
  \preceq\widehat P_r-I
  \preceq\frac{r}{1-r}I,
\]
and
\begin{equation}
  \norm{\widetilde T}
  \leq\frac{1}{\sqrt{r(1-r)}}.
  \label{eq:r-bound}
\end{equation}
The product \(r(1-r)\) is uniquely maximized on \((0,1)\) at
\(r=1/2\).  Substituting that value into \eqref{eq:r-bound} gives the
best bound produced by this family:
\[
  \norm{\widetilde T}\leq2.
\]

\subsection*{The canonical orbit Gramian and a controlled similarity}

The raw matrix \(P_r\) depends on the chosen eigenbasis, but the
balanced Gramian has a canonical, unitarily covariant form.  Since
\(T=U\widetilde T U^*\), define
\begin{equation}
  \mathcal P_r(T)
  =
  \sum_{k=0}^{\infty}r^{2k}T^{*k}T^k
  =
  U\widehat P_rU^*.
  \label{eq:canonical-Pr}
\end{equation}
It satisfies
\begin{equation}
  I\preceq\mathcal P_r(T)\preceq\frac1{1-r}I,
  \qquad
  \mathcal P_r(T)
  =
  I+r^2T^*\mathcal P_r(T)T.
  \label{eq:canonical-Pr-bounds}
\end{equation}
Put \(X_r=\mathcal P_r(T)^{1/2}\).  The second identity gives
\[
  (X_rTX_r^{-1})^*(X_rTX_r^{-1})
  =
  \frac1{r^2}
  \left(I-\mathcal P_r(T)^{-1}\right)
  \preceq\frac1r I,
\]
while the first gives
\[
  \kappa(X_r)
  =\norm{X_r}\norm{X_r^{-1}}
  \leq\frac1{\sqrt{1-r}}.
\]
Thus \(T\) is similar, through a quantitatively controlled positive
metric, to an operator of norm at most \(r^{-1/2}\).  Multiplying the
similarity bound and the transformed norm recovers
\[
  \norm{T}
  \leq
  \frac1{\sqrt{1-r}}\frac1{\sqrt r}.
\]
At \(r=1/2\), both factors equal \(\sqrt2\).  The constant \(2\) is
therefore the balanced product of a metric-conditioning cost and a
norm-in-that-metric cost.

\begin{intuitionbox}[Workshop extension]
The parameter calculation in this section is a pedagogical extension
of the source's half-eigenvalue sampling.  It shows that \(1/2\) is not
a lucky guess within this one-parameter family: it optimally balances
how strongly the detector sees the orbit against how strongly the
Stein recursion amplifies the norm.
\end{intuitionbox}

\section{Why the same-side algebra is insufficient}

Return to the matrices in \cref{ex:wrong-orientation}:
\[
  T_M=\begin{pmatrix}0&M\\0&1/2\end{pmatrix},
  \qquad H(w)\equiv I,
  \qquad M>0.
\]
They have distinct eigenvalues in \(\D\), and \(H\) is normalized
positive-real.  Moreover,
\[
  H(w)-(I-wT_M)^{-1}
  =-\sum_{k=1}^{\infty}w^kT_M^k\in\alg(T_M).
\]
Yet \(\norm{T_M}\geq M\).  Thus a same-side algebra constraint gives
no uniform bound.  The adjoint orientation in
\(\alg(T^*)\) is what becomes row-diagonal after the positivity-preserving
congruence.

\section{The reusable new-tool template}

The proof suggests the following five-step design pattern.
\begin{proofmap}[Structured kernel cancellation--Gramian rigidity]
\begin{enumerate}
  \item \textbf{Find a positive kernel} generated by a normalized
  positive-real analytic function.
  \item \textbf{Identify an adjoint-aligned nuisance module} that has
  fewer effective directions than an arbitrary matrix correction.
  \item \textbf{Build an annihilating sample}---often by adding an
  anchor point whose vector solves a linear system.
  \item \textbf{Recognize the surviving terms as ordered weighted
  Gramians}; their difference should be positive and detect the
  dynamics.
  \item \textbf{Use a Stein recursion} to turn the Gramian order bound
  into a sharp operator-norm estimate.
\end{enumerate}
\end{proofmap}
The method is credible only when every stage can be verified without
estimating the nuisance term.

\section{Core-proof oral checklist}

A participant is ready to proceed when they can answer, without notes:
\begin{enumerate}
  \item Why does simple spectrum make the defect diagonal after the
  chosen congruence?
  \item Why is \(\Theta(0)=0\) essential?
  \item Where do the denominators \(1-x/4\) and \(1-x/2\) come from?
  \item Why does \(v=-G^{-1}Pu\) cancel every \(\theta_i(w_i)\)?
  \item Why is \(\widehat Y\succeq0\)?
  \item Why does \(\widehat Ye=0\) imply
  \(\widehat Pe=e\)?
  \item Which identity finally turns \(2I-\widehat P\succeq0\) into
  \(\norm{T}\leq2\)?
\end{enumerate}

\section{Exercises}

\begin{exercise}[Warm-up: a scalar detector]
Let \(P=\alpha\geq1\), \(X=\beta\geq0\), and
\(cX-XP-PX\geq0\).  Determine what follows when \(\beta>0\) and what
additional information is needed when \(\beta=0\).
\end{exercise}

\begin{exercise}[Core: rigidity]
Prove \cref{lem:detector-rigidity} from scratch, including the use of
\cref{lem:psd-zero}.
\end{exercise}

\begin{exercise}[Core: Stein shift]
\label{ex:stein-shift}
Derive \eqref{eq:stein} directly from the series and use it to reproduce
the last operator inequality.
\end{exercise}

\begin{exercise}[Core: general radius]
\label{ex:general-radius}
Verify \eqref{eq:general-partial-fraction}, derive
\eqref{eq:general-anticommutator}, and optimize
\eqref{eq:r-bound}.
\end{exercise}

\begin{exercise}[Core: wrong algebra]
Complete \cref{ex:wrong-orientation} and explain why commutativity of
\(\alg(T)\) does not rescue the theorem.
\end{exercise}

\begin{exercise}[Stretch: equality clues]
Without consulting \cref{rem:auxiliary-strict}, reconstruct the
equality analysis as follows.  Assume \(\norm{\widetilde T}=2\) and
choose a unit vector \(e\) with \(\norm{\widetilde Te}=2\).  Trace
equality through the final chain to show
\[
  \widehat Pe=2e,
  \qquad
  \widehat P\widetilde Te=\widetilde Te.
\]
Deduce \(\widetilde T^2e=0\), and explain why this contradicts the
simple-spectrum hypothesis.  Then compare with the limiting nilpotent
restriction in \cref{cor:nilpotent-rigidity}.  Finally, examine the
diagonalizable matrices
\[
  T_\varepsilon=
  \begin{bmatrix}
    \varepsilon&2\\
    0&-\varepsilon
  \end{bmatrix},
  \qquad \varepsilon\neq0,
\]
and explain how \(T_\varepsilon\to
\begin{bmatrix}0&2\\0&0\end{bmatrix}\) illustrates the loss of
diagonalizability at an extremal limit.
\end{exercise}

\part{Removing hypotheses and completing the theorem}

\chapter{Perturbation and smooth convex outer approximation}
\label{ch:approximation}

\begin{learninggoals}
\begin{itemize}
  \item Approximate arbitrary matrices by simple-spectrum matrices.
  \item Force algebra generation by perturbing the scalar function.
  \item Construct \(C^1\) convex outer domains.
  \item Interpret parallel-body smoothing as a uniform support-slack
  regularization and derive an averaged resolvent bound.
  \item Keep a precise ledger of which limit removes which hypothesis.
\end{itemize}
\end{learninggoals}

\begin{workshopplan}
\textbf{Opening:} list every auxiliary hypothesis still present.
\textbf{Core I:} simple spectrum and the separating perturbation
\(f_\eta\). \textbf{Core II:} parallel bodies and \(C^1\) boundary.
\textbf{Studio:} write an explicit three-limit ledger.
\end{workshopplan}

\section{Simple-spectrum matrices are dense}

\begin{lemma}[Density of simple spectrum]
\label{lem:simple-density}
Matrices with \(n\) distinct eigenvalues are dense in \(\Mn\).
\end{lemma}

\begin{proof}
The discriminant of a monic polynomial vanishes exactly when the
polynomial has a repeated root.  Applied to the characteristic
polynomial, it is a polynomial in the entries of the matrix.  It is
not the zero polynomial: at \(\diag(1,2,\ldots,n)\) it is nonzero.
The zero set of a nonzero complex polynomial has empty interior.
Therefore every neighborhood in \(\Mn\) contains a matrix whose
characteristic discriminant is nonzero, hence whose eigenvalues are
distinct.
\end{proof}

\begin{blackbox}[Why a polynomial zero set has no interior]
For one complex variable, a nonzero polynomial has only finitely many
zeros.  For several variables, induct on the number of variables:
write the polynomial as a polynomial in the last variable whose
coefficients are polynomials in the others.  If it vanished on an open
polydisk, the one-variable statement would force every coefficient to
vanish on an open set; induction would make every coefficient the zero
polynomial.  This contradicts the assumption.  Applied to the real and
imaginary matrix entries, this is the density fact used above.
\end{blackbox}

Combine this with \cref{lem:W-Lipschitz}.  If
\(W(A)\) is compactly contained in an open domain \(\Omega\), set
\[
  \delta_\Omega=\dist(W(A),\C\setminus\Omega)>0.
\]
Whenever \(\norm{B-A}<\delta_\Omega\),
\[
  W(B)\subset W(A)+\norm{B-A}\olD\subset\Omega.
\]
Thus we may choose simple-spectrum \(B_k\to A\) while keeping all
their numerical ranges in the same domain.

\section{Why \texorpdfstring{\(f(B)\)}{f(B)} may not generate
\texorpdfstring{\(\alg(B)\)}{alg(B)}}

Let \(B\) have distinct eigenvalues
\(\beta_1,\ldots,\beta_n\).  If \(f(\beta_i)=f(\beta_j)\) for some
\(i\neq j\), then \(T=f(B)\) has repeated eigenvalues and cannot
recover the two corresponding spectral coordinates.  Both hypotheses
needed by the completion theorem can fail at once:
\[
  T\text{ may not have simple spectrum},\qquad
  \alg(T)\text{ may be smaller than }\alg(B).
\]
The remedy separates the values of \(f\) without breaking the unit
bound.

\section{The algebra-generating perturbation}

Fix a bounded convex domain \(\Omega\), and suppose
\(\max_{\overline\Omega}|f|\leq1\).  Put
\[
  R_\Omega=\max_{z\in\overline\Omega}|z|.
\]
For \(\eta\in\C\), define
\begin{equation}
  f_\eta(z)=
  \frac{f(z)+\eta z}{1+|\eta|R_\Omega}.
  \label{eq:feta}
\end{equation}
For each fixed \(\eta\), this is holomorphic in \(z\), and
\[
  |f_\eta(z)|
  \leq\frac{1+|\eta|R_\Omega}
             {1+|\eta|R_\Omega}
  =1
  \qquad(z\in\overline\Omega).
\]
The appearance of \(|\eta|\) affects dependence on the parameter
\(\eta\), but not holomorphicity in the variable \(z\).

For \(i\neq j\), the collision equation
\[
  f(\beta_i)+\eta\beta_i
  =f(\beta_j)+\eta\beta_j
\]
excludes at most one value of \(\eta\), because
\(\beta_i\neq\beta_j\).  There are finitely many pairs, hence finitely
many exceptional parameters.  We can choose nonexceptional
\(\eta_\ell\to0\).

Set
\[
  T_\ell=f_{\eta_\ell}(B).
\]
Its eigenvalues
\(
f_{\eta_\ell}(\beta_i)
\)
are distinct, so \(T_\ell\) has simple spectrum.  Also
\(T_\ell\in\alg(B)\).  Conversely, Lagrange interpolation gives a
polynomial \(q_\ell\) such that
\[
  q_\ell(f_{\eta_\ell}(\beta_i))=\beta_i
  \qquad(1\leq i\leq n).
\]
Because \(B\) is diagonalizable,
\[
  q_\ell(T_\ell)=B.
\]
Therefore
\begin{equation}
  \alg(T_\ell)=\alg(B).
  \label{eq:algebra-equality-perturbed}
\end{equation}

\begin{pitfall}
Adding \(\eta z\) without the denominator in \eqref{eq:feta} separates
eigenvalue images but may violate \(|f_\eta|\leq1\).  The denominator
is what keeps the Cayley transform positive-real.
\end{pitfall}

\section{Fixed-domain estimate}

\begin{proposition}[Estimate on one smooth outer domain]
\label{prop:fixed-domain}
Let \(A\in\Mn\), let \(p\) be a polynomial, and let \(\Omega\) be a
bounded convex \(C^1\) domain with \(W(A)\subset\Omega\).  Then
\begin{equation}
  \norm{p(A)}
  \leq2\max_{z\in\overline\Omega}|p(z)|.
  \label{eq:fixed-domain}
\end{equation}
\end{proposition}

\begin{proof}
First suppose \(B\) has simple spectrum and \(W(B)\subset\Omega\).
Set
\[
  m_\Omega=\max_{\overline\Omega}|p|.
\]
If \(m_\Omega=0\), the polynomial vanishes on the open set
\(\Omega\), hence is identically zero.  Otherwise let
\(f=p/m_\Omega\), so \(\max_{\overline\Omega}|f|\leq1\).

Use the nonexceptional perturbations above.  Each
\(T_\ell=f_{\eta_\ell}(B)\) has simple spectrum and generates the
same algebra as \(B\).  The auxiliary sharp bound
\cref{cor:auxiliary} gives
\[
  \norm{f_{\eta_\ell}(B)}\leq2.
\]
As \(\eta_\ell\to0\), \(f_{\eta_\ell}(B)\to f(B)\).  Therefore
\[
  \norm{p(B)}=m_\Omega\norm{f(B)}\leq2m_\Omega.
\]

Now choose simple-spectrum \(B_k\to A\) with
\(W(B_k)\subset\Omega\).  Polynomial evaluation and the operator norm
are continuous, so passing \(k\to\infty\) proves
\eqref{eq:fixed-domain}.
\end{proof}

\section{Smooth convex outer approximation}

\begin{lemma}[Parallel bodies]
\label{lem:outer-approximation}
Let \(K\subset\C\) be nonempty, compact, and convex.  For
\(\delta>0\), define
\[
  \Omega_\delta=\set{z:\dist(z,K)<\delta}.
\]
Then:
\begin{enumerate}
  \item \(\Omega_\delta\) is a bounded convex domain containing \(K\);
  \item \(\partial\Omega_\delta\) is \(C^1\);
  \item \(\overline{\Omega_\delta}=K+\delta\olD\), and
  \[
    d_{\mathrm H}(\overline{\Omega_\delta},K)=\delta.
  \]
\end{enumerate}
\end{lemma}

\begin{proof}
The identity
\(\overline{\Omega_\delta}=K+\delta\olD\) gives boundedness,
convexity, and containment.  It gives Hausdorff distance at most
\(\delta\); equality follows by taking a supporting point of \(K\) in
any fixed direction and moving it outward by \(\delta\) in that
direction.

We prove the boundary regularity.  Every \(z\in\C\) has a unique
nearest point \(\pi_K(z)\in K\).  Existence follows from compactness.
For uniqueness, if two different minimizers existed, their midpoint
would be in \(K\) and the parallelogram identity would make it strictly
closer.  To see continuity rather than assume it, let
\(p=\pi_K(z)\) and \(q=\pi_K(w)\).  Minimality along the line segments
from \(p\) to \(q\) and from \(q\) to \(p\) gives the variational
inequalities
\[
  \RePart\bigl(\overline{z-p}(q-p)\bigr)\leq0,
  \qquad
  \RePart\bigl(\overline{w-q}(p-q)\bigr)\leq0.
\]
Adding them yields
\(
|p-q|^2\leq\RePart(\overline{z-w}(p-q))
\),
and Cauchy--Schwarz gives
\[
  |\pi_K(z)-\pi_K(w)|\leq|z-w|.
\]
Thus the metric projection is nonexpansive, hence continuous.

Let \(d_K(z)=\dist(z,K)\).  Comparing \(z+h\) first with
\(\pi_K(z)\), and \(z\) with \(\pi_K(z+h)\), gives
\begin{align*}
2\RePart\bigl(
  \overline{z+h-\pi_K(z+h)}\,h
\bigr)-|h|^2
&\leq d_K(z+h)^2-d_K(z)^2\\
&\leq
2\RePart\bigl(
  \overline{z-\pi_K(z)}\,h
\bigr)+|h|^2.
\end{align*}
Continuity of \(\pi_K\) shows that \(d_K^2\) is continuously
differentiable, with real gradient
\begin{equation}
  \nabla d_K^2(z)=2(z-\pi_K(z)).
  \label{eq:distance-gradient}
\end{equation}
On the level set \(d_K(z)=\delta\), this gradient has length
\(2\delta\), so it never vanishes.  Since
\[
  \partial\Omega_\delta
  =\set{z:d_K(z)^2=\delta^2},
\]
the implicit-function theorem gives a \(C^1\) boundary.
\end{proof}

\begin{intuitionbox}
Adding a small disk rounds every corner just enough for a continuous
normal, without destroying convexity.  The squared distance function
is the defining equation, and its nonzero gradient is the outward
normal direction.
\end{intuitionbox}

\begin{proposition}[Parallel-body boundary resolvent bound]
\label{prop:parallel-resolvent-bound}
Let \(B\in\Mn\), put \(K=W(B)\), and let
\[
  \Omega_\delta=\set{z:\dist(z,K)<\delta},
  \qquad \delta>0.
\]
For \(\zeta\in\partial\Omega_\delta\), write
\(\mathcal R_\zeta=(\zeta I-B)^{-1}\).  Then
\begin{equation}
  \int_{\partial\Omega_\delta}
  \mathcal R_\zeta^*\mathcal R_\zeta\,ds
  \preceq\frac{2\pi}{\delta}I.
  \label{eq:parallel-resolvent-bound}
\end{equation}
Equivalently, for every \(x\in\C^n\),
\[
  \int_{\partial\Omega_\delta}
  \norm{\mathcal R_\zeta x}_2^2\,ds
  \leq\frac{2\pi}{\delta}\norm{x}_2^2.
\]
The constant is sharp already when \(B=0\).
\end{proposition}

\begin{proof}
Fix \(\zeta\in\partial\Omega_\delta\), put
\(p=\pi_K(\zeta)\), and use \eqref{eq:distance-gradient} to identify
the outward normal:
\[
  \nu(\zeta)=\frac{\zeta-p}{\delta}.
\]
The nearest-point variational inequality says that \(p\) is a support
point of \(K\) in the direction \(\nu(\zeta)\).  Hence
\[
  \RePart(\overline{\nu(\zeta)}p)
  =
  \lambda_{\max}
  \left(\RePart(\overline{\nu(\zeta)}B)\right).
\]
Since \(\zeta=p+\delta\nu(\zeta)\), the support slack from
\cref{prop:support-slack-test} is
\begin{align*}
  D_{\Omega_\delta,B}(\nu(\zeta))
  &=
  \delta I+
  \lambda_{\max}
  \left(\RePart(\overline{\nu(\zeta)}B)\right)I
  -\RePart(\overline{\nu(\zeta)}B)\\
  &\succeq\delta I.
\end{align*}
Apply \cref{prop:support-resolvent-resolution}:
\[
  2\pi I
  =
  \int_{\partial\Omega_\delta}
  \mathcal R_\zeta^*
  D_{\Omega_\delta,B}(\nu(\zeta))
  \mathcal R_\zeta\,ds
  \succeq
  \delta
  \int_{\partial\Omega_\delta}
  \mathcal R_\zeta^*\mathcal R_\zeta\,ds.
\]
For \(B=0\), direct integration on
\(\partial\Omega_\delta=\set{|\zeta|=\delta}\) gives equality.
\end{proof}

\begin{intuitionbox}[Geometric smoothing is operator regularization]
The support function of \(W(B)+\delta\olD\) is the support function of
\(W(B)\) plus \(\delta\).  Thus adding the disk adds exactly
\(\delta I\) to every directional support slack.  The weighted
conservation law from Week 8 then becomes the unweighted
vector-valued resolvent estimate
\eqref{eq:parallel-resolvent-bound}.
\end{intuitionbox}

\begin{pitfall}
Equation \eqref{eq:parallel-resolvent-bound} controls the integral of
\(\mathcal R_\zeta^*\mathcal R_\zeta\) in operator order, equivalently
one fixed vector at a time.  It does not imply the dimension-free
scalar inequality
\(\int\!\norm{\mathcal R_\zeta}^2ds\leq2\pi/\delta\).
\end{pitfall}

\section{The limit ledger}

There are three conceptually distinct limits:
\begin{equation}
  \boxed{
  \eta_\ell\to0
  \quad\longrightarrow\quad
  B_k\to A
  \quad\longrightarrow\quad
  \overline{\Omega_{\delta_j}}\to W(A).
  }
  \label{eq:limit-ledger}
\end{equation}
\begin{enumerate}
  \item \(\eta_\ell\to0\) removes simple spectrum and algebra
  generation for \(f(B)\), while \(B\) itself still has simple
  spectrum.
  \item \(B_k\to A\) removes simple spectrum for the matrix.
  \item \(\delta_j\downarrow0\) removes the smooth outer domain and
  returns to the exact numerical range.
\end{enumerate}
Writing ``by approximation'' without this order hides which theorem is
available at each stage.

\section{Exercises}

\begin{exercise}[Warm-up: separate two values]
Given distinct \(\beta_1,\beta_2\) but equal values
\(f(\beta_1)=f(\beta_2)\), find all \(\eta\) for which
\(f(\beta_j)+\eta\beta_j\) still collide.
\end{exercise}

\begin{exercise}[Core: exceptional parameters]
\label{ex:exceptional-eta}
For a simple-spectrum \(3\times3\) diagonal \(B\), write every
exceptional value of \(\eta\) explicitly in terms of
\(\beta_i\) and \(f(\beta_i)\).
\end{exercise}

\begin{exercise}[Core: algebra generation]
Prove \eqref{eq:algebra-equality-perturbed} in both directions, naming
the polynomial used for each inclusion.
\end{exercise}

\begin{exercise}[Core: fixed domain]
Reconstruct \cref{prop:fixed-domain} and mark the exact line where
each auxiliary hypothesis is used.
\end{exercise}

\begin{exercise}[Core: projection gradient]
\label{ex:projection-gradient}
Derive the two-sided inequality in the proof of
\cref{lem:outer-approximation} and then prove
\eqref{eq:distance-gradient}.
\end{exercise}

\begin{exercise}[Stretch: nonexpansive projection]
Prove
\[
  |\pi_K(z)-\pi_K(w)|\leq|z-w|
\]
using the variational inequalities for the two nearest points.
\end{exercise}

\chapter{Assembly of the main theorem and the spectral-set conclusion}
\label{ch:main-assembly}

\begin{learninggoals}
\begin{itemize}
  \item Reproduce the main polynomial proof in dependency order.
  \item Pass from outer-domain maxima to the numerical-range maximum.
  \item Pass the stronger orbit-Gramian certificate through all three
  approximation layers.
  \item Derive exact and quantitative rigidity for extremal pairs.
  \item Extend the polynomial estimate to rational functions.
  \item Distinguish theorem, classical input, and limiting step.
\end{itemize}
\end{learninggoals}

\begin{workshopplan}
\textbf{Opening:} place the ten proof dependencies in order.
\textbf{Core:} shrink smooth outer domains and pass the polynomial
bound to \(W(A)\). \textbf{Guided proof:} use a pole-free neighborhood
and Runge approximation for rational functions.
\textbf{Studio:} submit the full proof write-up with every temporary
hypothesis crossed out where it is removed.
\end{workshopplan}

\section{The polynomial estimate}

\begin{proof}[Proof of \cref{thm:main}, polynomial part]
Fix \(A\in\Mn\) and \(p\in\C[z]\).
The numerical range \(K=W(A)\) is nonempty, compact, and convex.
For \(j\geq1\), let
\[
  \Omega_j=\set{z:\dist(z,K)<1/j}.
\]
By \cref{lem:outer-approximation}, these are bounded convex \(C^1\)
domains containing \(K\), and
\[
  d_{\mathrm H}(\overline\Omega_j,K)\to0.
\]
Apply the fixed-domain estimate:
\[
  \norm{p(A)}
  \leq2\max_{z\in\overline\Omega_j}|p(z)|.
\]
All the sets eventually lie in one compact neighborhood of \(K\).
Uniform continuity of \(|p|\), or the Hausdorff-maxima exercise from
Week 5, gives
\[
  \max_{z\in\overline\Omega_j}|p(z)|
  \longrightarrow
  \max_{z\in W(A)}|p(z)|.
\]
Pass to the limit:
\[
  \norm{p(A)}
  \leq2\max_{z\in W(A)}|p(z)|.
\]
\end{proof}

\section{A stronger orbit-Gramian conclusion}

The approximation argument preserves more than the first-power norm
bound.  The canonical orbit Gramian from
\eqref{eq:canonical-Pr} passes through every limit even when the
diagonalizing eigenbases themselves become singular.

\begin{corollary}[Global orbit-Gramian estimate]
\label{cor:global-orbit-gramian}
Let \(p\in\C[z]\), and put
\[
  m=\max_{z\in W(A)}|p(z)|.
\]
If \(m>0\), define \(T=p(A)/m\).  Then, for every \(0<r<1\),
\begin{equation}
  \mathcal P_r(T)
  :=
  \sum_{k=0}^{\infty}r^{2k}T^{*k}T^k
  \preceq\frac1{1-r}I.
  \label{eq:global-orbit-gramian}
\end{equation}
In particular, at \(r=1/2\),
\begin{equation}
  \sum_{k=0}^{\infty}4^{-k}
  \norm{T^kx}_2^2
  \leq2\norm{x}_2^2
  \qquad(x\in\C^n).
  \label{eq:half-orbit-square-function}
\end{equation}
\end{corollary}

\begin{proof}
Fix \(0<r<1\).  At the auxiliary simple-spectrum,
algebra-generating stage, Week 12
proved
\[
  \mathcal P_r(X)\preceq\frac1{1-r}I.
\]
We record why this order bound survives the three limits in
\eqref{eq:limit-ledger}.  For fixed \(N\), put
\[
  S_N(X)=\sum_{k=0}^{N}r^{2k}X^{*k}X^k.
\]
If \(X_j\to X\) and
\(\mathcal P_r(X_j)\preceq(1-r)^{-1}I\), then polynomial continuity
gives
\[
  S_N(X)
  =
  \lim_{j\to\infty}S_N(X_j)
  \preceq\frac1{1-r}I.
\]
Apply this observation successively as the scalar perturbation tends
to zero, the simple-spectrum matrices tend to \(A\), and the outer
domains shrink to \(W(A)\).  The matrices at the final outer-domain
stage are
\[
  \frac{p(A)}
       {\max_{\overline\Omega_j}|p|}
  \longrightarrow T.
\]
Thus every finite partial sum \(S_N(T)\) is bounded above by
\((1-r)^{-1}I\).  The partial sums increase in positive-semidefinite
order, so they converge in finite dimensions and give
\eqref{eq:global-orbit-gramian}.  Equation
\eqref{eq:half-orbit-square-function} is its quadratic-form version
at \(r=1/2\).
\end{proof}

\begin{corollary}[Global controlled similarity]
\label{cor:global-controlled-similarity}
With the notation of \cref{cor:global-orbit-gramian}, put
\(X_r=\mathcal P_r(T)^{1/2}\).  Then
\[
  \kappa(X_r)\leq\frac1{\sqrt{1-r}},
  \qquad
  \norm{X_rTX_r^{-1}}\leq\frac1{\sqrt r}.
\]
In particular, at \(r=1/2\), the normalized matrix \(T\) is similar,
with condition number at most \(\sqrt2\), to a matrix of norm at most
\(\sqrt2\).
\end{corollary}

\begin{proof}
Shifting the convergent series gives the Stein identity
\[
  \mathcal P_r(T)
  =
  I+r^2T^*\mathcal P_r(T)T.
\]
Now repeat the calculation following
\eqref{eq:canonical-Pr-bounds}.  The order bounds
\[
  I\preceq\mathcal P_r(T)\preceq(1-r)^{-1}I
\]
control both the transformed norm and the condition number.
\end{proof}

\begin{remark}
The raw Cauchy matrix \(P_r\) depends on a diagonalizing eigenbasis and
need not converge when eigenvalues coalesce.  The canonical object
\(\mathcal P_r(T)\) is basis-free; passing finite partial sums is what
makes the strengthened conclusion stable.  If \(m=0\), then
the polynomial estimate just proved gives \(p(A)=0\), so there is no
normalization to perform.
\end{remark}

\begin{corollary}[Nilpotent rigidity of exact equality]
\label{cor:nilpotent-rigidity}
Under the hypotheses of \cref{cor:global-orbit-gramian}, suppose
\[
  \norm{p(A)}=2m.
\]
Then there are orthonormal vectors \(x,y\in\C^n\) such that
\[
  Tx=2y,
  \qquad
  Ty=0.
\]
Consequently, \(\Span\{y,x\}\) is \(T\)-invariant and, in the ordered
basis \((y,x)\),
\begin{equation}
  T|_{\Span\{y,x\}}
  =
  \begin{pmatrix}0&2\\0&0\end{pmatrix}.
  \label{eq:extremal-nilpotent-restriction}
\end{equation}
In particular, \(p(A)\), and hence \(A\), is not diagonalizable.
\end{corollary}

\begin{proof}
Choose a unit vector \(x\) with \(\norm{Tx}_2=2\), and let
\[
  P=\mathcal P_{1/2}(T).
\]
By \cref{cor:global-orbit-gramian},
\[
  2
  \geq x^*Px
  =
  1+\frac14\norm{Tx}_2^2
  +\sum_{k=2}^{\infty}4^{-k}\norm{T^kx}_2^2
  \geq2.
\]
Thus \(x^*(2I-P)x=0\).  Since \(2I-P\succeq0\),
\cref{lem:psd-zero} gives \(Px=2x\); the nonnegative tail also gives
\(T^2x=0\).  Put \(y=Tx/2\).  Then \(y\) is a unit vector and \(Ty=0\),
so the Gramian series gives \(Py=y\).  Since \(P\) is Hermitian, its
eigenvectors \(x\) and \(y\), belonging to the distinct eigenvalues
\(2\) and \(1\), are orthogonal.  This proves
\eqref{eq:extremal-nilpotent-restriction}.

The relations \(T^2x=0\) and \(Tx\neq0\) show that \(T\) is not
diagonalizable.  A polynomial of a diagonalizable matrix is
diagonalizable, so \(A\) cannot be diagonalizable either.  Also
\(0\in\spec(p(A))\), hence \(p\) vanishes at some eigenvalue of \(A\).
\end{proof}

\begin{pitfall}
The two-dimensional subspace in
\eqref{eq:extremal-nilpotent-restriction} is invariant for \(T\), but
need not reduce \(T\) and need not be invariant for \(A\).  The
corollary is a necessary condition for equality, not a classification
of all extremal pairs.
\end{pitfall}

\begin{corollary}[Quantitative near-extremal chain]
\label{cor:near-extremal-chain}
Let \(T=p(A)/m\) as above, and choose a unit vector \(x\) with
\(\norm{Tx}_2=\norm{T}\).  Then
\begin{equation}
  \norm{T^2x}_2^2
  \leq
  4\left(4-\norm{T}^2\right).
  \label{eq:near-extremal-chain}
\end{equation}
\end{corollary}

\begin{proof}
The \(k=0\) and \(k=1\) terms in
\eqref{eq:half-orbit-square-function} give
\[
  \sum_{k=2}^{\infty}4^{-k}\norm{T^kx}_2^2
  \leq
  1-\frac14\norm{T}^2
  =
  \frac{4-\norm{T}^2}{4}.
\]
Retain only the \(k=2\) term and multiply by \(16\).
\end{proof}

\section{The proof in dependency order}

\clearpage
\begin{proofmap}[Dependency-ordered reconstruction]
\begin{enumerate}
  \item Continuity on the finite-dimensional unit sphere makes
  \(W(A)\) compact; Toeplitz--Hausdorff makes it convex.
  \item A smooth convex outer domain supplies a positive double-layer
  map \(\Phi\).
  \item The Cayley family gives a normalized positive-real \(H\) whose
  resolvent defect lies in \(\alg(T^*)\), provided
  \(T=f(B)\) generates \(\alg(B)\).
  \item Simple spectrum diagonalizes that algebraic defect.
  \item Half-eigenvalue samples plus one origin sample cancel it.
  \item The surviving kernel inequality compares two weighted
  Gramians.
  \item Detector rigidity and the Stein identity yield
  \(\norm{T}\leq2\).
  \item The \(f_\eta\) perturbation enforces the two auxiliary
  hypotheses and then disappears.
  \item Simple-spectrum matrices approach \(A\) inside the fixed
  domain.
  \item Parallel outer domains shrink to \(W(A)\).
\end{enumerate}
\end{proofmap}

\section{From polynomials to rational functions}

Let \(r\) be rational with no pole on \(W(A)\).  Its finite pole set is
disjoint from the compact numerical range, so choose \(\delta>0\) such
that the compact convex neighborhood
\[
  K_\delta=W(A)+\delta\olD
\]
contains no pole.  Then \(r\) is holomorphic on a neighborhood of
\(K_\delta\), and \(\C\setminus K_\delta\) is connected.  Runge
approximation gives polynomials \(p_j\) with
\[
  p_j\longrightarrow r
  \quad\text{uniformly on }K_\delta.
\]
Choose a contour
\(\Gamma\subset\operatorname{int}K_\delta\) enclosing
\(\spec(A)\).  Uniform convergence on \(\Gamma\), together with the
Cauchy functional calculus, gives
\[
  p_j(A)\longrightarrow r(A).
\]
Uniform convergence on \(W(A)\) also gives
\[
  \max_{W(A)}|p_j|
  \longrightarrow\max_{W(A)}|r|.
\]
Apply the polynomial estimate and pass to the limit:
\[
  \norm{r(A)}
  \leq2\max_{z\in W(A)}|r(z)|.
\]
Thus \(W(A)\) is a \(2\)-spectral set for \(A\).

\begin{pitfall}
Uniform approximation only on \(W(A)\) is not the stated reason that
\(p_j(A)\to r(A)\).  The proof approximates on a pole-free
\emph{neighborhood} and uses convergence on a contour around the
spectrum.
\end{pitfall}

\section{Proof-audit table}

\begin{longtable}{@{}L{0.28\textwidth}L{0.31\textwidth}
                         L{0.32\textwidth}@{}}
\toprule
\textbf{Temporary assumption}&\textbf{Why it appears}&
\textbf{How it is removed}\\
\midrule
\endhead
\(\partial\Omega\) is \(C^1\)&Defines normal and double-layer density&
Parallel bodies of \(W(A)\)\\
\(W(A)\subset\Omega\)&Keeps resolvents and boundary calculus valid&
Hausdorff shrinkage\\
\(B\) has simple spectrum&Diagonal defect and Lagrange coordinates&
Density plus \(B_k\to A\)\\
\(f_\eta(B)\) has simple spectrum&Completion theorem hypothesis&
Finite exceptional set, then \(\eta\to0\)\\
\(\alg(f_\eta(B))=\alg(B)\)&Places companion defect in
\(\alg(T^*)\)&Same separating perturbation\\
\(|f_\eta|\leq1\)&Cayley positivity&
Normalization denominator\\
\bottomrule
\end{longtable}

\section{Exercises}

\begin{exercise}[Warm-up: shrinking disks]
For \(K=\olD\), write \(\Omega_j=K+(1/j)\D\) explicitly and verify
directly that
\(\max_{\overline\Omega_j}|z|^m\to\max_K|z|^m\).
\end{exercise}

\begin{exercise}[Core: two-page proof]
\label{ex:two-page-proof}
Write a self-contained two-page proof of the polynomial estimate,
citing but not reproving the positive-real completion theorem and the
double-layer lemma.
\end{exercise}

\begin{exercise}[Core: maximum convergence]
Give an \(\eps\)-proof that
\[
  \max_{\overline\Omega_j}|p|
  \to\max_{W(A)}|p|.
\]
\end{exercise}

\begin{exercise}[Core: rational passage]
\label{ex:rational-passage}
Reproduce the rational approximation argument and explain why
\(\spec(A)\subset W(A)\) is used.
\end{exercise}

\begin{exercise}[Stretch: replace \(p\)]
Show that the same argument proves
\[
  \norm{f(A)}
  \leq2\max_{W(A)}|f|
\]
for every function \(f\) holomorphic on a neighborhood of \(W(A)\).
\end{exercise}

\chapter{Consequences, computation, and applications}
\label{ch:consequences}

\begin{learninggoals}
\begin{itemize}
  \item Apply the theorem to holomorphic matrix functions.
  \item Prove the Hilbert-space polynomial form by compression.
  \item Turn scalar approximation error into matrix approximation error.
  \item Design a numerical-range experiment without confusing evidence
  with proof.
\end{itemize}
\end{learninggoals}

\begin{workshopplan}
\textbf{Opening:} turn a scalar approximation error into a matrix
error. \textbf{Core:} holomorphic and Hilbert-space consequences.
\textbf{Computation lab:} sample support lines of a numerical range
and calibrate on the exact nilpotent disk.
\textbf{Closing:} separate numerical evidence from proof.
\end{workshopplan}

\section{Holomorphic functions}

\begin{corollary}[Holomorphic form]
\label{cor:holomorphic}
If \(f\) is holomorphic on a neighborhood of \(W(A)\), then
\[
  \norm{f(A)}
  \leq2\max_{z\in W(A)}|f(z)|.
\]
\end{corollary}

\begin{proof}
Choose a compact convex neighborhood of \(W(A)\) on which \(f\) is
holomorphic.  Runge approximation gives polynomials converging
uniformly there.  Matrix evaluation converges by the Cauchy functional
calculus, and the numerical-range maxima converge uniformly.
\end{proof}

If \(s\) is a polynomial, or a rational function without poles on
\(W(A)\), apply the corollary to \(f-s\):
\begin{equation}
  \norm{f(A)-s(A)}
  \leq2\max_{z\in W(A)}|f(z)-s(z)|.
  \label{eq:matrix-error}
\end{equation}
This says that a good scalar approximant on the numerical range is,
up to the sharp universal factor \(2\), a good matrix approximant.

\section{Finite-dimensional compression on a Hilbert space}

Let \(\mathcal H\neq\{0\}\) be a complex Hilbert space, with inner
product conjugate-linear in the first variable, and let
\[
  W(A)=\set{\ip{x}{Ax}:\norm{x}=1}
\]
for \(A\in\mathcal B(\mathcal H)\).

\begin{corollary}[Hilbert-space polynomial form]
\label{cor:hilbert}
For every polynomial \(p\),
\[
  \norm{p(A)}
  \leq2\sup_{z\in W(A)}|p(z)|.
\]
Consequently, \(\overline{W(A)}\) is a \(2\)-spectral set for \(A\).
\end{corollary}

\begin{proof}
If \(p=0\), the claim is immediate.  Otherwise let \(d=\deg p\) and
fix a unit vector \(x\).  Define the finite
Krylov space
\[
  \mathcal M_x=\Span\set{x,Ax,\ldots,A^dx}
\]
and the compression
\[
  A_x=P_{\mathcal M_x}A|_{\mathcal M_x}.
\]
Inductively,
\begin{equation}
  A_x^kx=A^kx
  \qquad(0\leq k\leq d),
  \label{eq:krylov-agreement}
\end{equation}
because \(A^{k+1}x\in\mathcal M_x\) while \(k<d\).
Also \(W(A_x)\subset W(A)\).  The finite-dimensional theorem gives
\[
  \norm{p(A)x}
  =\norm{p(A_x)x}
  \leq\norm{p(A_x)}
  \leq2\sup_{z\in W(A)}|p(z)|.
\]
Take the supremum over unit \(x\).

For the rational conclusion we use the standard inclusion
\begin{equation}
  \spec(A)\subseteq\overline{W(A)}
  \label{eq:hilbert-spectrum-in-W}
\end{equation}
for a bounded operator on a complex Hilbert space.  If a rational
function \(r\) has no pole on \(\overline{W(A)}\), compactness of its
finite pole set gives a pole-free compact neighborhood of
\(\overline{W(A)}\).  Approximate \(r\) there by polynomials and use a
contour around \(\spec(A)\) inside that neighborhood.  The polynomial
bound and the Cauchy functional calculus then pass to the limit exactly
as in Week 14.
\end{proof}

\begin{blackbox}[Spectrum lies in the closed numerical range]
If \(\lambda\notin\overline{W(A)}\), a separating line and a rotation
give \(c>0\) such that
\[
  \RePart\ip{x}{e^{-i\theta}(A-\lambda I)x}\geq c\norm{x}^2
\]
for every \(x\).  Cauchy--Schwarz shows that \(A-\lambda I\) is bounded
below.  The analogous inequality for its adjoint shows that the
adjoint has trivial kernel, so the range is both closed and dense,
hence all of the Hilbert space.  Therefore \(A-\lambda I\) is
invertible, proving \eqref{eq:hilbert-spectrum-in-W}.
\end{blackbox}

\section{A numerical-range plotting algorithm}

Formula \eqref{eq:W-support-function} suggests a robust boundary
sampler:
\begin{enumerate}
  \item choose angles \(\theta_k=2\pi k/N\);
  \item form the Hermitian matrix
  \[
    H_{\theta_k}=\RePart(e^{-i\theta_k}A);
  \]
  \item compute a normalized eigenvector \(x_k\) for its largest
  eigenvalue;
  \item record \(z_k=x_k^*Ax_k\).
\end{enumerate}
The line with outward normal \(e^{i\theta_k}\) supports \(W(A)\) at
\(z_k\).  Repeated largest eigenvalues may produce flat faces, so one
should sample the compression of
\(\ImPart(e^{-i\theta_k}A)\) to the top eigenspace if the entire face
is desired.

\section{A workshop experiment}

For selected matrices and polynomials, approximate the Crouzeix ratio
\[
  \mathcal C(A,p)=
  \frac{\norm{p(A)}}
       {\max_{z\in W(A)}|p(z)|}.
\]
A careful experiment should include:
\begin{itemize}
  \item a normal matrix, where the ratio is at most \(1\);
  \item the nilpotent sharp example, where the ratio equals \(2\);
  \item a mildly nonnormal triangular matrix;
  \item a random \(3\times3\) matrix and several polynomial degrees;
  \item mesh-refinement checks for the denominator.
\end{itemize}
Numerical evidence can test intuition and reveal near-extremizers, but
it cannot certify the universal inequality.

\section{Why approximation bounds matter}

Many matrix algorithms approximate \(f(A)b\) by a polynomial or
rational expression \(s(A)b\).  Equation \eqref{eq:matrix-error}
separates the work:
\[
  \text{scalar approximation on }W(A)
  \quad\Longrightarrow\quad
  \text{matrix error control}.
\]
This perspective is relevant to polynomial and rational Krylov
methods, including GMRES; see the survey
\cite{BickelEtAl2020} and the numerical-range spectral-set literature.

\section{Exercises}

\begin{exercise}[Core: holomorphic form]
Prove \cref{cor:holomorphic} with an explicit surrounding compact set
and contour.
\end{exercise}

\begin{exercise}[Core: Krylov agreement]
\label{ex:krylov}
Prove \eqref{eq:krylov-agreement} by induction and verify
\(W(A_x)\subset W(A)\).
\end{exercise}

\begin{exercise}[Core: error transfer]
Derive \eqref{eq:matrix-error} for both polynomial and rational \(s\).
\end{exercise}

\begin{exercise}[Project: numerical-range atlas]
Implement the support-line sampler.  Report the matrix, the sampled
boundary, at least three polynomials, the observed ratios, and a
mesh-refinement study.  Include the exact nilpotent calculation as a
calibration.
\end{exercise}

\begin{exercise}[Stretch: flat faces]
Choose a concrete \(3\times3\) matrix for which the largest eigenvalue
of \(H_\theta\) has multiplicity two.  Compute the tangential
compression \(K_\theta\) from
\cref{prop:exposed-face-pencil} and use its extreme eigenvalues to
decide whether the supported face is a point or a nontrivial segment.
Give one example of each behavior.
\end{exercise}

\chapter{Scope, open directions, and a new-tool clinic}
\label{ch:scope}

\begin{learninggoals}
\begin{itemize}
  \item Distinguish the scalar theorem from complete spectral-set
  questions.
  \item Separate the operator-valued support geometry of \(W(A)\) from
  the domain-independent completion mechanism.
  \item Identify the hypotheses that make kernel cancellation work.
  \item Formulate plausible extensions without claiming them as
  theorems.
  \item Present the proof and its new tool coherently.
\end{itemize}
\end{learninggoals}

\begin{workshopplan}
\textbf{Opening:} distinguish scalar from complete spectral-set
questions. \textbf{Core:} diagnose which hypotheses power the new
tool and which are artifacts of this proof.
\textbf{Capstone studio:} present one proof stage, one failure mode,
and one plausible extension, labeling conjectures as conjectures.
\end{workshopplan}

\section{What has been established}

The course has developed the source manuscript's scalar statement:
\[
  \norm{p(A)}
  \leq2\max_{z\in W(A)}|p(z)|.
\]
It controls one scalar-valued polynomial at a time and yields the
rational \(2\)-spectral-set conclusion.

\section{What is specifically special about the numerical range}

The proof uses no curvature formula, conformal parametrization, or
strict convexity of \(W(A)\).  For a bounded \(C^1\) convex test
domain \(\Omega\), define its support function on all of \(\Torus\) by
\[
  h_\Omega(u)
  =
  \max_{z\in\overline\Omega}\RePart(\overline uz).
\]
When \(\partial\Omega\cap\spec(A)=\varnothing\), the
numerical-range-specific bridge is
\[
  W(A)\subseteq\overline\Omega
  \quad\Longleftrightarrow\quad
  h_\Omega(u)I-\RePart(\overline uA)\succeq0
  \quad(u\in\Torus),
\]
which resolvent congruence turns into the positive boundary density
\eqref{eq:support-slack-factorization}.  In other words, every scalar
supporting-half-plane inequality lifts exactly to operator order.
The numerical range is the intersection, and the Hausdorff limit, of
the smooth convex outer enclosures producing these positive kernels.
This formulation allows \(W(A)\) itself to have corners or to meet the
spectrum on its boundary, where the boundary resolvent need not be
defined.

After this bridge has produced the positive-real completion, the sharp
argument is independent of the planar shape.  The factor \(2\) comes
from the optimally balanced sampling scale \(r=1/2\), the ordered
Gramians, and the Stein identity.  Thus the useful structure belongs
to the pair consisting of the operator and its numerical range, not
to the unlabeled convex set alone.

\section{What is stronger and remains separate}

For a matrix polynomial
\[
  \mathcal F(z)=[f_{ij}(z)]\in\mathbb M_m(\C[z]),
\]
one can ask whether
\begin{equation}
  \norm{\mathcal F(A)}
  \leq2\max_{z\in W(A)}\norm{\mathcal F(z)}
  \label{eq:matrix-level-question}
\end{equation}
holds uniformly in \(m\).  Here \(\mathcal F(A)\) acts on
\((\C^n)^m\).  The case \(m=1\) is the scalar theorem.  Uniformity in
\(m\) is an additional operator-algebraic property and does not follow
formally from the scalar case; compare
\cite{DavidsonPaulsenWoerdeman2018}.

In the block-valued setting, a diagonal scalar nuisance becomes a
family of noncommuting blocks.  The single equation \(Gv+Pu=0\) no
longer obviously annihilates every nuisance direction.  A block or
bimodule version of the cancellation principle would be needed.

Notice that the geometry-to-boundary-map step is already completely
positive by the POVM construction in Week 8.  The matrix-level
obstruction is therefore not a failure of convex positivity.  It is
the algebraic de-symmetrization step: at matrix level the correction
no longer reduces to scalar diagonal entries that one origin vector
can cancel.

\section{Potential extensions of the abstract tool}

The positive-real completion theorem does not mention numerical
ranges.  This suggests several research programs:
\begin{itemize}
  \item replace a finite diagonal algebra by a semisimple commutative
  algebra or a simultaneously diagonalizable tuple;
  \item find tangential or block annihilators for matrix-valued
  functions;
  \item seek forward--backward Gramian pairs for annular domains;
  \item replace a finite eigenbasis by spectral measures or a Riesz
  basis in infinite dimensions;
  \item classify equality beyond the necessary nilpotent restriction
  in \cref{cor:nilpotent-rigidity}, and sharpen the stability estimate
  \eqref{eq:near-extremal-chain}.
\end{itemize}
These are directions, not consequences proved in this text.

\section{Capstone options}

\begin{enumerate}
  \item Give a line-by-line expository proof of
  \cref{thm:completion}, with a symbolic \(2\times2\) audit.
  \item Derive the general-radius estimate
  \eqref{eq:r-bound} and explain its optimization geometrically.
  \item Compute the double-layer measure on a disk and compare it with
  the Poisson kernel.
  \item Starting from \cref{cor:nilpotent-rigidity}, investigate which
  additional conditions make the invariant nilpotent restriction
  sufficient for equality.
  \item Compare scalar spectral sets with the matrix-level question
  \eqref{eq:matrix-level-question}.
  \item Design an abstract nuisance module and determine whether a
  finite anchor sample can annihilate it.
\end{enumerate}

\section{Final conceptual summary}

\begin{proofmap}[The idea to remember]
The earlier symmetrized boundary identity contains a useful companion
term but does not isolate \(f(B)\).  The source proof retains the full
Cayley family, revealing that the resolvent defect lies in an
adjoint-generated algebra.  Simple spectrum turns that algebraic fact
into a diagonal analytic correction.  Herglotz-kernel samples at half
the conjugate eigenvalues, together with a compensating sample at the
origin, eliminate the correction exactly.  The remaining inequality
orders two weighted Gramians.  Their positive difference detects the
dynamics, and the narrower Gramian satisfies a Stein identity.  That
combination forces the sharp norm bound \(2\).

The proof also leaves behind more structure than the norm estimate:
the support slacks resolve the identity through the boundary
resolvents, the canonical orbit Gramians survive all approximation
layers, and exact saturation forces the normalized functional-calculus
output
\[
  T=\frac{p(A)}{\max_{z\in W(A)}\abs{p(z)}}
\]
to contain an invariant \(2\times2\) nilpotent restriction.
\end{proofmap}

\section{Research-status note}

The repository describes the source as a \emph{candidate proof} for
which formal peer review is pending.  It also records computational,
adversarial, and formalization audits that found no specific error.
Accordingly, this workshop presents the manuscript's argument in a
checkable form and preserves the distinction between a research
manuscript, formal verification in a stated system, and independent
mathematical peer review.  The additional support-slack, orbit-Gramian,
and equality-rigidity results in these notes are presented as
deductions from the manuscript's mechanism, not as claims of priority
in the literature.

\section{Exercises}

\begin{exercise}[Core: scalar versus complete]
Explain why applying the scalar theorem to each entry of
\(\mathcal F\) does not prove \eqref{eq:matrix-level-question} with
constant \(2\).
\end{exercise}

\begin{exercise}[Core: tool diagnosis]
For a proposed extension, identify: the positive kernel, the nuisance
module, the annihilator equation, the two ordered Gramians, and the
Stein recursion.  If one component is missing, state the obstruction.
\end{exercise}

\begin{exercise}[Capstone oral]
Give a fifteen-minute proof overview in which every classical input,
new step, and approximation step is explicitly labeled.
\end{exercise}

\begin{exercise}[Stretch: stability]
Suppose the final norm is close to \(2\).  Trace which inequalities
must be nearly equalities and formulate a quantitative stability
question.
\end{exercise}

\appendix

\chapter{Proof dependencies and equation ledger}
\label{app:dependencies}

\section{Minimal dependency graph}

The main proof can be checked in four independent packets.

\subsection*{Packet A: finite-dimensional linear algebra}

\begin{itemize}
  \item simple spectrum implies diagonalizability;
  \item Lagrange interpolation identifies generated algebras;
  \item congruence preserves positivity;
  \item the polar decomposition balances a nonorthogonal eigenbasis;
  \item Hermitian eigenvectors test matrix inequalities.
\end{itemize}

\subsection*{Packet B: analytic positivity}

\begin{itemize}
  \item scalar Herglotz representation;
  \item polarization to a positive matrix-valued measure;
  \item positivity of the matrix Herglotz kernel;
  \item scalar Cayley transform maps the disk to the right half-plane.
\end{itemize}

\subsection*{Packet C: numerical-range boundary calculus}

\begin{itemize}
  \item Toeplitz--Hausdorff convexity and
  \(\spec(B)\subset W(B)\);
  \item supporting-line geometry and its operator-valued support
  slacks;
  \item the positive double-layer density, its exact containment test,
  and mass \(2I\);
  \item the support--resolvent resolution of the identity and the
  parallel-body \(L^2\) resolvent bound;
  \item the boundary POVM and symmetrized normal dilation;
  \item the companion identity;
  \item smooth parallel outer domains.
\end{itemize}

\subsection*{Packet D: the new mechanism}

\begin{itemize}
  \item the Cayley family makes a resolvent completion;
  \item the defect lies in \(\alg(T^*)\);
  \item simple spectrum makes the defect diagonal after congruence;
  \item half-eigenvalue and origin samples cancel it;
  \item two Gramians, detector rigidity, and a Stein identity give
  \(2\);
  \item arbitrary-radius sampling gives a family of canonical
  orbit-Gramian certificates;
  \item those certificates survive approximation and force nilpotent
  rigidity at equality.
\end{itemize}

Packets A--C are classical ingredients.  Packet D is the
manuscript's central new mechanism.

\section{Normalization ledger}

\begin{longtable}{@{}L{0.17\textwidth}L{0.30\textwidth}
                         L{0.42\textwidth}@{}}
\toprule
\textbf{Quantity}&\textbf{Where it comes from}&\textbf{What it does}\\
\midrule
\endhead
\(2I\)&Total mass of two boundary layers&
Makes \(\Phi=\frac12\int d\mu_B\) unital\\
\(2\)&Cayley expansion \(1+2\sum_{m\geq1}w^mf^m\)&
Turns the known double-layer coefficient into \(T^m\)\\
\(1/2\)&Sample \(w_i=\overline{\lambda_i}/2\)&
Produces the paired weights \(1/2\) and \(1/4\)\\
\(P\)&Kernel denominator \(1-w_i\overline{w_j}\)&
Determines the compensator \(v=-G^{-1}Pu\)\\
\(Q\)&Resolvent denominator \(1-w_i\lambda_j\)&
Creates the wider Gramian\\
\(4Q-2P\)&Partial fractions in the sample block&
Leads to the coefficient \(4\) in the anticommutator\\
\(4\)&Detector--anticommutator coefficient&
Gives \(\widehat P\preceq2I\)\\
\(1/4\)&Stein recursion weight&
Turns \(\widehat P\preceq2I\) into
\(\widetilde T^*\widetilde T\preceq4I\)\\
\bottomrule
\end{longtable}

\section{One-page completion proof}

Assume \(T=S\Lambda S^{-1}\), \(G=S^*S\), and
\[
  \widetilde H(w)
  =G(I-w\Lambda)^{-1}+\Theta(w)G,
  \qquad
  \Theta(0)=0,\quad\Theta(w)\text{ diagonal}.
\]
Set
\[
  P_{ij}=\frac{G_{ij}}{1-\overline{\lambda_i}\lambda_j/4},
  \quad
  Q_{ij}=\frac{G_{ij}}{1-\overline{\lambda_i}\lambda_j/2},
  \quad
  Y=Q-P.
\]
Sample the positive Herglotz kernel at
\[
  w_i=\overline{\lambda_i}/2,\quad
  \xi_i=u_ie_i,
  \qquad
  w_0=0,\quad
  \xi_0=-G^{-1}Pu.
\]
The correction contribution contains \(Gv+Pu=0\), hence vanishes.
The known blocks are \(4Q-2P\), \(G+Q\), and \(2G\), so
\[
  4Y-YG^{-1}P-PG^{-1}Y\succeq0.
\]
After balancing,
\[
  4\widehat Y-\widehat Y\widehat P
  -\widehat P\widehat Y\succeq0,
\]
with
\[
  \widehat P=\sum_{k\geq0}4^{-k}\widetilde T^{*k}\widetilde T^k,
  \quad
  \widehat Y=\sum_{k\geq1}(2^{-k}-4^{-k})
  \widetilde T^{*k}\widetilde T^k\succeq0.
\]
The detector lemma gives
\(
I\preceq\widehat P\preceq2I
\).
The Stein identity
\[
  \widehat P
  =I+\frac14\widetilde T^*\widehat P\widetilde T
\]
then gives
\[
  \widetilde T^*\widetilde T
  \preceq4(\widehat P-I)\preceq4I.
\]
Because \(T\) and \(\widetilde T\) are unitarily equivalent,
\(\norm{T}\leq2\).

\chapter{Hints and selected solutions}
\label{app:solutions}

\section{Weeks 1--4}

\subsection*{The sharp matrix norm (\cref{ex:norm-N})}

\[
  N^*N
  =
  \begin{pmatrix}0&0\\2&0\end{pmatrix}
  \begin{pmatrix}0&2\\0&0\end{pmatrix}
  =
  \begin{pmatrix}0&0\\0&4\end{pmatrix}.
\]
Its largest eigenvalue is \(4\), so \(\norm{N}=2\).

\subsection*{The \(2\times2\) PSD test (\cref{ex:psd-tests})}

Necessity of \(a,d\geq0\) follows by testing the coordinate vectors.
If \(a>0\), complete the square:
\[
  \begin{pmatrix}\overline x&\overline y\end{pmatrix}
  H\begin{pmatrix}x\\y\end{pmatrix}
  =
  a\left|x+\frac{b}{a}y\right|^2
  +\left(d-\frac{|b|^2}{a}\right)|y|^2.
\]
This is nonnegative exactly when \(ad\geq|b|^2\).  If \(a=0\), the
determinant condition forces \(b=0\), and \(d\geq0\) finishes.

\subsection*{Interpolation (\cref{ex:lagrange})}

For distinct \(\lambda_j\), set
\[
  \ell_j(z)=\prod_{k\neq j}
  \frac{z-\lambda_k}{\lambda_j-\lambda_k},
  \qquad
  q(z)=\sum_jd_j\ell_j(z).
\]
Then \(q(\lambda_j)=d_j\).  Substituting a diagonalization of \(T\)
gives \(q(T)=S\diag(d_j)S^{-1}\).

\subsection*{Equality of generated algebras (\cref{ex:same-algebra})}

\(f(B)\in\alg(B)\) gives
\(\alg(f(B))\subseteq\alg(B)\).  If the \(f(\beta_i)\) are distinct,
choose a polynomial \(q\) with \(q(f(\beta_i))=\beta_i\).  Since \(B\)
is diagonalizable, \(q(f(B))=B\), so the reverse inclusion holds.

\subsection*{The nilpotent disk (\cref{ex:nilpotent-disk})}

Up to a global phase, every unit vector is
\(
x=(\cos t,e^{i\phi}\sin t)^T
\).
Then
\[
  x^*Nx=e^{i\phi}\sin(2t).
\]
The radius \(\sin(2t)\) takes every value in \([0,1]\), and \(\phi\)
takes every angle.  Hence the image is exactly \(\olD\).

\section{Weeks 5--8}

\subsection*{Matrix Cauchy mass (\cref{ex:matrix-cauchy-mass})}

If \(B=S\Lambda S^{-1}\), integrate
\[
  (\zeta I-B)^{-1}
  =S\diag((\zeta-\lambda_j)^{-1})S^{-1}
\]
entrywise.  Each diagonal contour integral is \(1\), so the result is
\(SIS^{-1}=I\).  For a general matrix, apply the holomorphic
functional calculus to the constant function \(1\), or use Jordan
form and note that higher-order pole terms integrate to zero.

\subsection*{Herglotz factorization (\cref{ex:herglotz-factor})}

Use \(\overline\zeta=\zeta^{-1}\) and
\[
  \overline{\frac{\zeta+z}{\zeta-z}}
  =\frac{1+\overline z\zeta}{1-\overline z\zeta}.
\]
Bring the two fractions to a common denominator.  The numerator
becomes \(2(1-w\overline z)\), which cancels the outer factor.

\subsection*{Double-layer density (\cref{ex:density-positive})}

Let \(R=Z^{-1}\).  Then
\[
  R^*(\nu Z^*+\overline\nu Z)R
  =
  \nu R^*Z^*R+\overline\nu R^*ZR
  =
  \nu R+\overline\nu R^*.
\]
The middle expression is positive by congruence.

\subsection*{Companion identity (\cref{ex:companion-identity})}

Insert the boundary formula for \(g_h(z)\) into its Cauchy functional
calculus on an interior contour \(\Gamma\).  Interchange the two finite
contour integrals.  For fixed boundary \(\zeta\), the inner integral is
\[
  \frac{1}{2\pi i}\int_\Gamma
  \frac{(zI-B)^{-1}}{\zeta-z}\,dz
  =(\zeta I-B)^{-1}.
\]
The contour remains strictly inside \(\Omega\), so no boundary
extension of \(g_h\) is required.

\section{Weeks 9--12}

\subsection*{Cayley coefficient audit (\cref{ex:coefficient-audit})}

\[
  H(w)
  =\Phi(1)+2\sum_{m\geq1}w^m\Phi(f^m)
  =I+\sum_{m\geq1}w^m
   \bigl(T^m+g_{f^m}(B)^*\bigr).
\]
The first equality uses bounded linearity of \(\Phi\); the second uses
the double-layer identity once for every \(m\).

\subsection*{Wrong algebra orientation (\cref{ex:wrong-orientation})}

Since \(\rho(wT_M)<1\),
\[
  I-(I-wT_M)^{-1}
  =-\sum_{k\geq1}w^kT_M^k\in\alg(T_M).
\]
But
\[
  \norm{T_M}\geq\norm{T_Me_2}_2
  =\sqrt{M^2+\frac14}\to\infty.
\]
Thus the conclusion fails with \(\alg(T)\).

\subsection*{Correction cancellation (\cref{ex:theta-cancel})}

The half of the kernel quadratic form containing \(E(w_i)\) is
\[
  \sum_{i=1}^n
  (u_ie_i)^*E(w_i)
  \left(
    v+\sum_{j=1}^n
    \frac{u_je_j}{1-w_i\overline{w_j}}
  \right).
\]
Since \(E(w_i)=\Theta(w_i)G\), this is the expression inside the real
part in \eqref{eq:theta-contribution}.  The adjoint half is its complex
conjugate.  The denominator creates \(P\), so the bracket is
\((Gv+Pu)_i=0\).

\subsection*{Known top-left block (\cref{ex:top-left-block})}

For \(F(w)=G(I-w\Lambda)^{-1}\),
\[
  F(w_i)_{ij}
  =\frac{G_{ij}}{1-\overline{\lambda_i}\lambda_j/2},
\]
and \(F(w_j)^*_{ij}\) is the same.  Divide their sum by
\(
1-w_i\overline{w_j}
=1-\overline{\lambda_i}\lambda_j/4
\)
and apply the stated partial fraction identity.

\subsection*{Gramian expansion (\cref{ex:gramian-expansion})}

\[
  P=\sum_{k\geq0}4^{-k}\Lambda^{*k}G\Lambda^k.
\]
Congruence by \(G^{-1/2}\) changes the \(k\)-th summand into
\[
  G^{-1/2}\Lambda^{*k}G\Lambda^kG^{-1/2}
  =\widetilde T^{*k}\widetilde T^k.
\]
The calculation for \(Q\) is identical with \(2^{-k}\).

\subsection*{Stein shift (\cref{ex:stein-shift})}

Remove the \(k=0\) term from \(\widehat P\), factor one
\(\widetilde T^*\) on the left and one \(\widetilde T\) on the right,
and reindex.  Then
\[
  \widetilde T^*\widetilde T
  \preceq\widetilde T^*\widehat P\widetilde T
  =4(\widehat P-I)\preceq4I.
\]

\subsection*{General radius (\cref{ex:general-radius})}

Solving
\[
  \frac{2}{(1-rx)(1-r^2x)}
  =\frac{a}{1-rx}+\frac{b}{1-r^2x}
\]
gives \(a=2/(1-r)\) and \(b=-2r/(1-r)\).
After the same graph substitution, pure \(P_r\)-terms cancel because
\(a(1-r)=2\).  Rigidity gives
\(\widehat P_r\preceq(1-r)^{-1}I\), and the Stein identity yields
\(
\norm{\widetilde T}^2\leq[r(1-r)]^{-1}
\).
Differentiate \(r(1-r)\), or complete the square:
\[
  r(1-r)=\frac14-(r-\tfrac12)^2.
\]

\section{Weeks 13--16}

\subsection*{Exceptional parameters (\cref{ex:exceptional-eta})}

For each \(i<j\), a collision occurs only at
\[
  \eta_{ij}
  =-\frac{f(\beta_i)-f(\beta_j)}
          {\beta_i-\beta_j}.
\]
The common positive denominator in \(f_\eta\) does not affect
collisions.

\subsection*{Projection gradient (\cref{ex:projection-gradient})}

For the upper bound compare \(z+h\) with the candidate
\(\pi_K(z)\):
\[
  d_K(z+h)^2
  \leq|z+h-\pi_K(z)|^2.
\]
Expand and subtract \(d_K(z)^2\).  For the lower bound compare \(z\)
with \(\pi_K(z+h)\) and rearrange.  Divide the resulting two-sided
estimate by \(|h|\); continuity of \(\pi_K\) leaves the real-linear
derivative
\[
  h\longmapsto2\RePart(
  \overline{z-\pi_K(z)}h).
\]

\subsection*{Two-page proof outline (\cref{ex:two-page-proof})}

Use exactly four paragraphs:
\begin{enumerate}
  \item fixed \(\Omega\), simple-spectrum approximation \(B_k\);
  \item normalize \(p\), perturb \(f_\eta\), and obtain algebra
  generation;
  \item invoke the double-layer completion and completion theorem,
  then take \(\eta\to0\) and \(B_k\to A\);
  \item take parallel outer domains to \(W(A)\).
\end{enumerate}
Any proof that changes this order should explicitly justify why the
required hypotheses are still available.

\subsection*{Rational passage (\cref{ex:rational-passage})}

Choose a pole-free compact convex neighborhood \(K_\delta\), use
Runge on that set, and choose a contour around
\(\spec(A)\).  The inclusion \(\spec(A)\subset W(A)\) ensures that
the contour can be placed inside the pole-free neighborhood and that
the rational functional calculus is defined.

\subsection*{Krylov compression (\cref{ex:krylov})}

Assume \(A_x^kx=A^kx\) for \(k<d\).  Then
\[
  A_x^{k+1}x
  =P_{\mathcal M_x}AA^kx
  =A^{k+1}x,
\]
because \(A^{k+1}x\in\mathcal M_x\).  For a unit
\(y\in\mathcal M_x\),
\[
  \ip{y}{A_xy}
  =\ip{y}{P_{\mathcal M_x}Ay}
  =\ip{y}{Ay},
\]
so \(W(A_x)\subset W(A)\).

\chapter{Glossary and notation}
\label{app:glossary}

\begin{longtable}{@{}L{0.23\textwidth}L{0.68\textwidth}@{}}
\toprule
\textbf{Symbol or term}&\textbf{Meaning}\\
\midrule
\endhead
\(\C,\D,\Torus\)&Complex plane, open unit disk, and unit circle\\
\(\Mn\)&\(n\times n\) complex matrices\\
\(A^*\)&Conjugate transpose or Hilbert-space adjoint\\
\(\norm{x}_2\)&Euclidean vector norm\\
\(\norm{A}\)&Induced Euclidean operator norm\\
\(\spec(A)\)&Spectrum, the set of eigenvalues in finite dimensions\\
\(W(A)\)&Numerical range
\(\{x^*Ax:\norm{x}_2=1\}\)\\
\(\RePart A\)&Hermitian part \((A+A^*)/2\)\\
\(H\succeq0\)&\(H\) is positive semidefinite\\
\(H\preceq K\)&\(K-H\succeq0\)\\
\(\alg(T)\)&Unital algebra of polynomials in \(T\)\\
\(C\)-spectral set&Compact \(K\supset\spec(A)\) satisfying
\(\norm{r(A)}\leq C\max_K|r|\)\\
resolvent&\((zI-A)^{-1}\) when \(z\notin\spec(A)\)\\
simple spectrum&\(n\) distinct eigenvalues for an \(n\times n\) matrix\\
congruence&Transformation \(H\mapsto S^*HS\), preserving positivity\\
similarity&Transformation \(A\mapsto S^{-1}AS\), preserving spectrum\\
Carath\'eodory function&Analytic \(H\) on \(\D\) with
\(\RePart H\succeq0\)\\
positive kernel&Kernel satisfying every finite sampled quadratic
inequality \eqref{eq:positive-kernel}\\
\(\Phi\)&Unital positive double-layer map\\
POVM&Positive operator-valued measure of total mass \(I\)\\
\(g_h\)&Holomorphic Cauchy companion of boundary data \(\overline h\)\\
\(D_{\Omega,B}(u)\)&Directional support-slack matrix
\(h_\Omega(u)I-\RePart(\overline uB)\)\\
\(G=S^*S\)&Gram matrix of a diagonalizing eigenbasis\\
\(P,Q\)&Cauchy-type matrices that become two weighted Gramians\\
\(Y=Q-P\)&Positive detector after balancing\\
\(\mathcal P_r(T)\)&Canonical orbit Gramian
\(\sum_{k\geq0}r^{2k}T^{*k}T^k\)\\
Stein identity&Discrete recursion \(P=I+aT^*PT\)\\
\(d_{\mathrm H}\)&Hausdorff distance between compact sets\\
\bottomrule
\end{longtable}

\chapter{Classical inputs: exact statements for review}
\label{app:classical-inputs}

\section{Toeplitz--Hausdorff}

For every bounded operator \(A\) on a complex Hilbert space, its
numerical range is convex.  In finite dimensions it is also compact.
The course proves the finite-dimensional case by \(2\times2\)
compression in Week 3.

\section{Scalar Herglotz}

An analytic scalar function on \(\D\) has nonnegative real part if and
only if it has the representation
\[
  f(w)=ic+\int_{\Torus}\frac{\zeta+w}{\zeta-w}\,d\mu(\zeta),
\]
where \(c\in\R\) and \(\mu\) is a finite positive measure.  The
representation is unique.  Week 7 derives the matrix form and the
kernel consequence.

\section{Naimark dilation for a POVM}

Let \(E\) be a positive operator-valued measure on a compact space
\(X\), acting on a Hilbert space \(\mathcal H\), and suppose
\(E(X)=I_{\mathcal H}\).  Then there are a Hilbert space
\(\mathcal K\), an isometry \(V:\mathcal H\to\mathcal K\), and a
projection-valued measure \(F\) on \(X\) such that
\[
  E(\mathcal E)=V^*F(\mathcal E)V
\]
for every Borel set \(\mathcal E\subseteq X\).  Equivalently, the
unital positive map
\[
  \varphi\longmapsto\int_X\varphi\,dE
\]
is the compression of a unital \(*\)-representation of \(C(X)\), and
is therefore completely positive.

A proof constructs \(\mathcal K\) by completing vector-valued simple
functions modulo the nullspace of the semidefinite form
\[
  [f,g]_E=\int_X f(x)^*\,dE(x)\,g(x).
\]
Multiplication by characteristic functions gives the projections
\(F(\mathcal E)\), and constant functions define the isometry \(V\).
Week 8 applies this construction with \(X=\partial\Omega\).

\section{Runge}

If \(K\subset\C\) is compact with connected complement, every function
holomorphic near \(K\) is uniformly approximable on \(K\) by
polynomials.  The course applies this only to compact convex sets.

\section{Implicit-function theorem}

If a \(C^1\) real-valued function on \(\R^2\) has nonzero gradient on a
level set, that level set is locally a \(C^1\) curve.  Week 13 applies
this to \(d_K^2(z)=\delta^2\).

\section{Finite-dimensional polar decomposition}

Every invertible complex matrix has a unique factorization
\[
  S=U(S^*S)^{1/2},
\]
where \(U\) is unitary.  This converts the nonunitary diagonalization
into the unitary equivalence
\(
T=U(G^{1/2}\Lambda G^{-1/2})U^*
\).

\backmatter

\chapter*{Provenance and attribution}
\addcontentsline{toc}{chapter}{Provenance and attribution}

These course notes are based on the local source manuscript
\emph{The numerical range is a \(2\)-spectral set} by Shanmu Jin.
The source manuscript states that OpenAI ChatGPT contributed the idea
of sampling the matrix Herglotz kernel at half the conjugate
eigenvalues and adding the origin sample that cancels the
adjoint-algebra correction, as well as assisting with exposition,
bibliography, and typesetting.  The source assigns responsibility for
the mathematical statements and citations to its human author.

This workshop edition expands the definitions, prerequisites, proof
calculations, exercises, approximation ledger, and the general-radius
interpretation.  It also records further deductions from the proof:
the support-slack boundary calculus, its resolvent conservation law,
the canonical orbit-Gramian estimate, and nilpotent rigidity at
equality.  These additions should be read together with the
research-status note in Week 16.

\begin{thebibliography}{99}
\addcontentsline{toc}{chapter}{Bibliography}

\bibitem{JinManuscript}
S.~Jin,
\emph{The numerical range is a \(2\)-spectral set},
candidate manuscript in the accompanying research repository, 2026.

\bibitem{BadeaCrouzeixDelyon2006}
C.~Badea, M.~Crouzeix, and B.~Delyon,
Convex domains and \(K\)-spectral sets,
\emph{Math. Z.} \textbf{252} (2006), 345--365.

\bibitem{BickelEtAl2020}
K.~Bickel, P.~Gorkin, A.~Greenbaum, T.~Ransford,
F.~L.~Schwenninger, and E.~Wegert,
Crouzeix's conjecture and related problems,
\emph{Comput. Methods Funct. Theory} \textbf{20} (2020), 701--728.

\bibitem{Crouzeix2004}
M.~Crouzeix,
Bounds for analytical functions of matrices,
\emph{Integral Equations Operator Theory} \textbf{48} (2004),
461--477.

\bibitem{Crouzeix2007}
M.~Crouzeix,
Numerical range and functional calculus in Hilbert space,
\emph{J. Funct. Anal.} \textbf{244} (2007), 668--690.

\bibitem{CrouzeixGreenbaum2019}
M.~Crouzeix and A.~Greenbaum,
Spectral sets: numerical range and beyond,
\emph{SIAM J. Matrix Anal. Appl.} \textbf{40} (2019), 1087--1101.

\bibitem{CrouzeixPalencia2017}
M.~Crouzeix and C.~Palencia,
The numerical range is a \((1+\sqrt2)\)-spectral set,
\emph{SIAM J. Matrix Anal. Appl.} \textbf{38} (2017), 649--655.

\bibitem{DavidsonPaulsenWoerdeman2018}
K.~R.~Davidson, V.~I.~Paulsen, and H.~J.~Woerdeman,
Complete spectral sets and numerical range,
\emph{Proc. Amer. Math. Soc.} \textbf{146} (2018), 1189--1195.

\bibitem{DelyonDelyon1999}
B.~Delyon and F.~Delyon,
Generalization of von Neumann's spectral sets and integral
representation of operators,
\emph{Bull. Soc. Math. France} \textbf{127} (1999), 25--41.

\bibitem{RansfordSchwenninger2018}
T.~Ransford and F.~L.~Schwenninger,
Remarks on the Crouzeix--Palencia proof that the numerical range is a
\((1+\sqrt2)\)-spectral set,
\emph{SIAM J. Matrix Anal. Appl.} \textbf{39} (2018), 342--345.

\end{thebibliography}

\end{document}
